\documentclass[11pt,a4paper]{article}
\input{glyphtounicode}
\pdfgentounicode=1

\usepackage[utf8]{inputenc}
\usepackage[T1]{fontenc}
\usepackage{lmodern}
\usepackage{amsmath,amssymb,amsthm}
\usepackage{booktabs}
\usepackage{array}
\usepackage{longtable}
\usepackage{geometry}
\usepackage{graphicx}
\usepackage{xcolor}
\usepackage{hyperref}

\geometry{margin=1in}
\emergencystretch=1.5em
\hypersetup{
    colorlinks=true,
    linkcolor=blue!70!black,
    citecolor=blue!70!black,
    urlcolor=blue!70!black,
    pdftitle={Discrete Fundamental Physics from a Causal Poset and the Unique Minimal Finite Group G},
    pdfauthor={Gemini 4 Argon, Holger Joest},
    pdfcreationdate={D:20261115010000}
}

\newtheorem{theorem}{Theorem}[section]
\newtheorem{lemma}[theorem]{Lemma}
\newtheorem{proposition}[theorem]{Proposition}
\newtheorem{corollary}[theorem]{Corollary}
\theoremstyle{definition}
\newtheorem{definition}[theorem]{Definition}
\newtheorem{axiom}[theorem]{Axiom}
\theoremstyle{remark}
\newtheorem{remark}[theorem]{Remark}

\newcommand{\R}{\mathbb{R}}
\newcommand{\C}{\mathbb{C}}
\newcommand{\quat}{\mathbb{H}}
\newcommand{\Z}{\mathbb{Z}}
\newcommand{\Q}{\mathbb{Q}}
\newcommand{\F}{\mathbb{F}}
\newcommand{\Tr}{\operatorname{Tr}}
\newcommand{\STr}{\operatorname{STr}}
\newcommand{\Irr}{\operatorname{Irr}}
\newcommand{\SU}{\operatorname{SU}}
\newcommand{\SO}{\operatorname{SO}}
\newcommand{\Spin}{\operatorname{Spin}}
\newcommand{\U}{\operatorname{U}}
\newcommand{\SL}{\operatorname{SL}}
\newcommand{\GL}{\operatorname{GL}}
\newcommand{\Gthree}{G_3}
\newcommand{\Gstar}{G_\star}

\title{\textbf{Discrete Fundamental Physics from a Causal Poset and the Unique Minimal Finite Group $G_\star = 2I \times \Sigma(216 \times 3) \times \Delta(27)_f \times \Z_3^{\mathrm{PQ}}$}}
\author{Gemini 4 Argon, Holger Joest}
\date{15. November 2026}

\begin{document}

\maketitle

\begin{abstract}
We formulate and test a strictly finitist synthesis of Causal Set Theory (Bombelli, Lee, Meyer, and Sorkin~\cite{Bombelli1987}) and Noncommutative Geometry (Chamseddine and Connes~\cite{ChamseddineConnes1997}) in which both continuous spacetime and the internal Standard Model spectral triple $(\mathcal{A}_F = \C \oplus \quat \oplus M_3(\C), \mathcal{H}_F, D_F)$ emerge from three discrete structures: (i) a locally finite causal poset $(P, \le_P)$ of elementary spacetime events, (ii) a single finite group $G$ whose real group algebra $\R[G]$ encodes Lorentz spinor, gauge, Higgs, and flavor degrees of freedom via its Real Wedderburn decomposition, and (iii) a total event poset $Z \subseteq P \times G$ governed at each vertex by a local ordered product conservation law $\prod_{i=1}^m \pi_G(u_{\sigma(i)}) = \prod_{j=1}^n \pi_G(v_{\tau(j)})$. Analyzing the ascending sequence of minimal candidate groups $G_1 \subset G_2 \subset \Gthree \subset \Gstar$ across 40 structural and phenomenological consistency tests, we prove a \emph{5-Point Design Theorem} demonstrating that all smaller candidates ($G_1 = \Sigma(216 \times 3)$, $G_2 = 2I \times \Sigma(216 \times 3)$, and ungraded $\Gthree = 2I \times \Sigma(216 \times 3) \times \Delta(27)_f$) are ruled out by exact representation-theoretic no-go theorems, selecting the minimal four-factor group
\[
\Gstar = \underbrace{2I \times \Sigma(216 \times 3) \times \Delta(27)_f}_{\Gthree\;\text{(order }2{,}099{,}520\text{)}} \times \Z_3^{\mathrm{PQ}} \qquad (\text{order } 6{,}298{,}560).
\]
From discrete group-algebraic invariants and selection rules, the Plancherel expansion and Wedderburn decomposition of $\R[\Gstar]$ select vertex coordination degree $d = 4$, generate unitary quantum phase evolution $\exp(-iHt/\hbar)$ with exact non-supersymmetric Bose--Fermi supertrace balance, enforce $N_{\mathrm{gen}} = 3$ and all-orders proton stability ($p \to e^+\pi^0 = 0$) while permitting $\Delta B = \Delta L = 3$ sphaleron baryogenesis, and accommodate $\bar{\theta}_{\mathrm{QCD}}^{\mathrm{tree}} = 0$ with non-vanishing CKM CP violation ($J_{\mathrm{CP}} = 2.93 \times 10^{-5}$) on the Hermitian subspace of $\R[\Delta(27)_f]$. Using $M_Z$ as the unit-defining scale, the algebra reproduces the qualitative structure and leading-order magnitudes of 35 fundamental observables across 162 orders of magnitude---including $\alpha_{\mathrm{em}}^{-1}(M_Z) = 127.941$ ($\alpha_{\mathrm{em}}^{-1}(0) = 137.026$ with standard low-energy QED running), $\sin^2\theta_W(M_Z) = 0.23111$, $\alpha_s(M_Z) = 0.1170$, $\alpha_G(m_p) = 5.838 \times 10^{-39}$, $m_{h^0} = 125.10\text{ GeV}$, the charged-lepton, quark, and Dirac neutrino mass spectra and CKM/PMNS mixing matrices, a $2.462\text{ TeV}$ weak-triplet WIMP dark matter candidate ($\Omega_{\mathrm{DM}} h^2 = 0.1190$) co-located with a heavy Type-II 2HDM quartet, the baryon asymmetry $\eta_B = 6.125 \times 10^{-10}$, and the dark energy scale $E_\Lambda = 2.299\text{ meV}$.
\end{abstract}

\tableofcontents

%=============================================================================
\section{Introduction and Foundational Axioms}
%=============================================================================

Modern fundamental physics rests on two pillars---General Relativity on a continuous 4-dimensional Lorentzian manifold and the Standard Model of particle physics over the continuous gauge group $\SU(3)_c \times \SU(2)_{W,L} \times \U(1)_Y$---yet leaves unexplained why spacetime has $d = 3+1$ dimensions, why the internal gauge group and chiral fermion representations take their specific forms, why there are precisely $N_{\mathrm{gen}} = 3$ fermion generations, and what fixes the $25+$ dimensionless coupling constants, mass ratios, mixing angles, and cosmological hierarchies spanning $162$ orders of magnitude from $\rho_\Lambda / M_{\mathrm{Pl}}^4 \sim 10^{-123}$ to the top Yukawa coupling $y_t \sim 1$.

In this paper, we investigate a strictly finitist foundation in which both continuous spacetime manifolds and continuous Lie groups are emergent infrared approximations to a discrete algebraic structure $(P, G, Z, \prec_Z)$ defined by three axioms:

\begin{axiom}[Causal Spacetime Poset]
\label{ax:poset}
Spacetime events form a locally finite partially ordered set $(P, \le_P)$ whose covering relations $x \prec_P y$ ($x <_P y$ with no $w \in P$ satisfying $x <_P w <_P y$) represent elementary directed causal links.
\end{axiom}

\begin{axiom}[Finite Group Fiber and Real Group Algebra]
\label{ax:group}
Internal gauge, Lorentz spinor, Higgs, and flavor quantum numbers are carried by a single finite group $G$. Physical states and gauge/Yukawa operators belong to the semisimple real group algebra
\begin{equation}
\label{eq:wedderburn}
\R[G] \;\cong\; \bigoplus_{\alpha \in \Irr_{\R}(G)} M_{n_\alpha}(\mathbb{D}_\alpha), \qquad \mathbb{D}_\alpha \in \{\R, \C, \quat\}, \qquad \sum_{\chi \in \Irr_{\C}(G)} d_\chi^2 = |G|,
\end{equation}
where the division algebra $\mathbb{D}_\alpha$ of each simple Wedderburn block is uniquely determined by the Frobenius--Schur indicator
\begin{equation}
\label{eq:fs_indicator}
\iota(\chi) \;\equiv\; \frac{1}{|G|} \sum_{g \in G} \chi(g^2) \;\in\; \{+1, 0, -1\},
\end{equation}
corresponding to strictly real ($\iota = +1$, $\mathbb{D}_\alpha = \R$), complex ($\iota = 0$, $\mathbb{D}_\alpha = \C$), and quaternionic/pseudoreal ($\iota = -1$, $\mathbb{D}_\alpha = \quat$) sectors.
\end{axiom}

\begin{axiom}[Total Event Poset and Local Ordered Product Conservation Law]
\label{ax:conservation}
The total event space is a subset $Z \subseteq P \times G$ equipped with projections $\pi_P : Z \to P$ and $\pi_G : Z \to G$. To avoid the intracell covering-relation triviality trap (Proposition~\ref{prop:triviality_trap}), the partial order on $Z$ is induced strictly by the spacetime poset:
\begin{equation}
\label{eq:induced_order}
u <_Z v \iff \pi_P(u) <_P \pi_P(v), \qquad u \prec_Z v \iff \pi_P(u) \prec_P \pi_P(v).
\end{equation}
At every non-boundary event $z \in Z$ with $m = |U(z)| \ge 1$ immediate predecessors $U(z) = \{u_1, \dots, u_m\}$ and $n = |V(z)| \ge 1$ immediate successors $V(z) = \{v_1, \dots, v_n\}$, causal propagation satisfies the \emph{local ordered product conservation law}:
\begin{equation}
\label{eq:product_law}
\exists\, \sigma \in S_m,\, \tau \in S_n : \qquad \prod_{i=1}^m \pi_G(u_{\sigma(i)}) \;=\; \prod_{j=1}^n \pi_G(v_{\tau(j)}) \quad \text{in } G.
\end{equation}
\end{axiom}

\begin{proposition}[Resolution of the Covering-Relation Triviality Trap]
\label{prop:triviality_trap}
If a partial order $\le_Z$ on $Z = P \times G$ allowed intracell ordering $(x, g) <_Z (x, g')$ for $g \neq g'$ on a dense fiber $Z_x$, then for any $x \prec_P y$, $(x, g_1)$ and $(y, g_2)$ would cease to be covering relations whenever a third element $(x, g_3)$ or $(y, g_3)$ intervened, disconnecting causal transport between distinct spacetime points. Enforcing Eq.~\eqref{eq:induced_order} makes every fiber $Z_x = \pi_P^{-1}(x)$ an antichain and guarantees well-defined causal covering neighborhoods $U(z)$ and $V(z)$.
\end{proposition}

\begin{theorem}[Abelianization Kirchhoff Conservation Law]
\label{thm:kirchhoff}
Let $\phi : G \to G^{\mathrm{ab}} = G / [G,G]$ be the canonical projection onto the abelianization of $G$. Because $G^{\mathrm{ab}}$ is abelian, all permutations $\sigma \in S_m$ and $\tau \in S_n$ collapse to the same sum, and Eq.~\eqref{eq:product_law} projects at every vertex $z \in Z$ to an exact, permutation-independent discrete Kirchhoff current conservation law:
\begin{equation}
\label{eq:kirchhoff}
\sum_{i=1}^m \phi(\pi_G(u_i)) \;=\; \sum_{j=1}^n \phi(\pi_G(v_j)) \quad \text{in } G^{\mathrm{ab}}.
\end{equation}
Consequently, for any balanced vertex coordination degree $m = n = d \to \infty$, the joint compatibility probability $p_{\mathrm{joint}}(d)$ of random group tuples is bounded above by $1/|G^{\mathrm{ab}}|$, providing a permanent non-washout selectivity floor of $\log_2|G^{\mathrm{ab}}|$ bits per vertex.
\end{theorem}

%=============================================================================
\section{Representation-Theoretic Classification and the 5-Point Design Theorem}
%=============================================================================

Rather than postulating a finite group $G$ ad hoc, we analyze the ascending hierarchy of minimal subgroup products $G_1 \subset G_2 \subset \Gthree \subset \Gstar$:
\begin{itemize}
    \item \textbf{Single-Factor Gauge Candidate ($G_1$):} $G_1 = \Sigma(216 \times 3)$ (the double cover of the Hessian group, order $648$).
    \item \textbf{Two-Factor Spin--Gauge Candidate ($G_2$):} $G_2 = 2I \times \Sigma(216 \times 3)$ (direct product with the binary icosahedral group $2I \subset \SU(2)$, order $120 \times 648 = 77{,}760$).
    \item \textbf{Three-Factor Spin--Gauge--Flavor Subgroup ($\Gthree$):} $\Gthree = 2I \times \Sigma(216 \times 3) \times \Delta(27)_f$ (direct product with the non-abelian Heisenberg flavor group $\Delta(27)_f \cong (\Z_3 \times \Z_3) \rtimes \Z_3$, order $120 \times 648 \times 27 = 2{,}099{,}520$, with abelianization $|\Gthree^{\mathrm{ab}}| = 27$).
    \item \textbf{Complete Four-Factor Group ($\Gstar$):} $\Gstar = \Gthree \times \Z_3^{\mathrm{PQ}} = 2I \times \Sigma(216 \times 3) \times \Delta(27)_f \times \Z_3^{\mathrm{PQ}}$ (order $6{,}298{,}560$).
\end{itemize}

\subsection{Exact No-Go Theorems Eliminating $G_1$, $G_2$, and Ungraded $\Gthree$}

\begin{theorem}[No-Go Theorem 1: 4D Crystallographic Periodicity of $G_1 = \Sigma(216 \times 3)$]
\label{thm:nogo1}
In $G_1 = \Sigma(216 \times 3) \cong \Delta(27)_c \rtimes 2T_W$, every irreducible representation $\chi$ of real dimension $\le 4$ ($\chi_0, \dots, \chi_6$) contains the $27$-element Heisenberg color subgroup $\Delta(27)_c$ in its kernel ($|\ker\chi| \ge 27$, image order $\le 24$) and possesses strictly integer real traces $\Tr_{\R}\rho(g) \in \{-4,-2,-1,0,1,2,3,4\} \subset \Z$. By the 4D Crystallographic Restriction Theorem, any 4D spacetime projection built from $G_1$ alone is strictly periodic (generating the $D_4/F_4$ root lattice of the $24$-cell) and cannot produce an aperiodic 4D quasicrystal.
\end{theorem}

\begin{theorem}[No-Go Theorem 2: Abelian Generation Collapse over $\R$ and Circulant CKM Triviality in $G_1$ and $G_2$]
\label{thm:nogo2}
Suppose the three fermion generations arise from tensoring a base representation $\rho_0$ with the three $1$-dimensional characters $\{1, \lambda, \lambda^2\}$ of an abelian $\Z_3$ factor (as in $G_1$ and $G_2$, where $\Sigma(216 \times 3)^{\mathrm{ab}} \cong \Z_3$). Then:
\begin{enumerate}
    \item \textbf{Lepton Generation Collapse:} For any color-singlet lepton representation $\rho_0$ that is real or quaternionic ($\overline{\rho}_0 \simeq \rho_0$), complex conjugation swaps the $\lambda$ and $\lambda^2$ generations ($\overline{\rho_0 \otimes \lambda} \simeq \rho_0 \otimes \lambda^2$). In $\R[G]$, the two non-trivial generations fuse into a single simple Wedderburn block, collapsing $3$ generations to $1\text{ real/quaternionic} + 1\text{ complex}$ generation and destroying quark--lepton universality.
    \item \textbf{Circulant CKM Identity $|V_{\mathrm{CKM}}| = I_3$:} Because $\bar{Q}_L^{(a)} \otimes q_R^{(b)} \simeq (\bar{Q}_L^{(0)} \otimes q_R^{(0)}) \otimes \lambda^{b-a}$ depends solely on the cyclic difference $(b - a) \bmod 3$, every $3 \times 3$ Yukawa matrix $M_u$ and $M_d$ is strictly $\Z_3$-circulant and simultaneously diagonalized by the discrete Fourier matrix $U_\omega = \frac{1}{\sqrt{3}}(\omega^{jk})$, forcing $V_{\mathrm{CKM}} = U_u^\dagger U_d$ to be a trivial phase-permutation matrix ($\theta_{12} = \theta_{23} = \theta_{13} = 0$, $J_{\mathrm{CP}} = 0$).
\end{enumerate}
\end{theorem}

\begin{theorem}[No-Go Theorem 3: The Conjugate-Higgs Reality Trap, Color-Breaking Octet VEVs, and the Two-Job Overload in Ungraded $\Gthree$]
\label{thm:nogo3}
In the ungraded three-factor group $\Gthree = 2I \times \Sigma(216 \times 3) \times \Delta(27)_f$, $\Sigma(216 \times 3)$ possesses only one $\Z_3$ abelianization ($2T_W^{\mathrm{ab}} \cong \Z_3$) and $\Delta(27)_f$ possesses only one $3$D irrep $\Delta[9]$ of central triality $+1$:
\begin{enumerate}
    \item \textbf{Conjugate-Higgs Reality Trap:} If $u_R \sim \Sigma[11](\lambda^{+1})$ and $d_R \sim \Sigma[10](\lambda^{-1})$ couple to the conjugate pair $\Sigma[4](\bar{2}_0) \oplus \Sigma[3](2_0) = M_2(\C)$, reality in $\R[\Gthree]$ forces $v_{p,q}^{(u)} = (v_{-p,-q}^{(d)})^*$ and $\Tr(M_u M_u^\dagger) \equiv \Tr(M_d M_d^\dagger)$ ($m_t \approx m_b$).
    \item \textbf{Color-Octet VEV Obstruction:} Attempting to split $m_t \neq m_b$ via a VEV in the shared real octet $\Sigma[19](8_{+1}) \in M_8(\R)$ is ruled out because $|\ker\Sigma[19]| = 3 < 27$: $\Sigma[19]|_{\Delta(27)_c}$ decomposes into the $8$ non-trivial color characters, so $\langle\Phi_{\Sigma[19]}\rangle \neq 0$ spontaneously breaks $\SU(3)_c$ color at the electroweak scale.
    \item \textbf{The Hypercharge vs.\ Peccei--Quinn Two-Job Overload:} Using $\Sigma(216 \times 3)^{\mathrm{ab}} \cong \Z_3$ simultaneously as weak hypercharge and as a Type-II 2HDM Peccei--Quinn grading ($u_R \sim \Sigma[7](\lambda^0)$, $e_R \sim \Sigma[0](\lambda^0)$) puts $e_R$ in the trivial representation ($Q_{\mathrm{em}}(e_R) = 0$) and leaves continuous hypercharge anomalous ($[\SU(3)_c]^2 \U(1)_Y \neq 0$ in $\R$).
\end{enumerate}
\end{theorem}

\subsection{The 5-Point Design Theorem and Minimality Proof for $\Gstar$}

\begin{theorem}[The 5-Point Design Theorem and Minimality of $\Gstar$]
\label{thm:design}
Any finite group $G$ on $(P, G, Z, \prec_Z)$ that reproduces a 4D aperiodic spacetime with chiral spinors, an unbroken $\SU(3)_c \times \SU(2)_{W,L} \times \U(1)_Y$ gauge sector with $\sin^2\theta_W(\Lambda) = 3/8$, three complex fermion generations over $\R[G]$ with non-trivial CKM/PMNS mixing, and $m_t \gg m_b$ without breaking $\SU(3)_c$ or neutralizing $e_R$ must contain four minimal factors:
\begin{enumerate}
    \item \textbf{Spacetime/Spinor Factor $G_{\mathrm{spin}} = 2I$ (Order 120):} Unique minimal subgroup of $\Spin(4) \cong \SU(2)_L \times \SU(2)_R$ with two faithful quaternionic $2$D spinor doublets $\mathbf{2}_{\mathrm{spin},L} \oplus \mathbf{2}_{\mathrm{spin},R}$ ($\iota = -1$) and $\Q(\sqrt{5})$ traces projecting $E_8 \to H_4 \oplus H_4'$.
    \item \textbf{Gauge Factor $G_{\mathrm{gauge}} = \Sigma(216 \times 3) \cong \Delta(27)_c \rtimes 2T_W$ (Order 648):} Since $\sum d_\chi^2 = |G|$, $\Delta(27)_c$ ($9 \cdot 1^2 + 2 \cdot 3^2 = 27$) is the strictly smallest group with a faithful $3$D complex color triplet ($\iota = 0$), $2T_W \cong \SL(2,\F_3)$ (order $24$) is the strictly smallest group with a faithful quaternionic weak doublet ($\iota = -1$) and complex hypercharge characters $\lambda^{\pm 1}$ ($\iota = 0$), and their semidirect product fuses the $8$ non-trivial color characters into the irreducible real gluon octet $\Sigma[19](\mathbf{8}_{+1}) \in M_8(\R)$.
    \item \textbf{Non-Abelian Flavor Factor $G_{\mathrm{fl}} = \Delta(27)_f$ (Order 27):} Strictly smallest group supplying a faithful $3$D complex generation triplet $\mathbf{3}_f = \Delta[9]$ ($\iota = 0$) whose bilinear $\overline{\mathbf{3}}_f \otimes \mathbf{3}_f = \bigoplus_{(p,q)\in\Z_3^2} \chi_{p,q}^{(f)}$ spans all $9$ non-commuting Weyl--Heisenberg operators $W_{p,q} = P^p Q^q \in M_3(\C)$.
    \item \textbf{Discrete Peccei--Quinn Factor $G_{\mathrm{PQ}} = \Z_3^{\mathrm{PQ}}$ (Order 3):} Supplies the Type-II 2HDM charges $q_{\mathrm{PQ}}(Q_L, u_R, d_R, L_L, e_R, \nu_R) = (0, 0, -1, +1, 0, +1)$ separating $H_u = (\Sigma[4], \mathrm{PQ}[0]) \in M_2(\C)_{q_{\mathrm{PQ}}=0}$, $H_d = H_e = (\Sigma[3], \mathrm{PQ}[-1]) \in M_2(\C)_{q_{\mathrm{PQ}}=\mp 1}$, and $H_\nu = (\Sigma[5], \mathrm{PQ}[0]) \in \quat_{q_{\mathrm{PQ}}=0}$ into three disjoint color-singlet Wedderburn blocks while forbidding all six Dimension-4 RPV operators.
    \item \textbf{The 3-Generation Discrete Anomaly Theorem:} A single generation with $q_{\mathrm{PQ}} = (0, 0, -1, +1, 0, +1)$ has discrete Ib\'a\~nez--Ross anomalies $(\mathcal{A}_{\SU(3)_c^2\Z_3}^{(1)}, \mathcal{A}_{\SU(2)_L^2\Z_3}^{(1)}, \mathcal{A}_{\mathrm{Grav}^2\Z_3}^{(1)}) \equiv (1, 1, 1) \pmod 3$. Total discrete anomaly cancellation holds if and only if $3 \mid N_{\mathrm{gen}}$, which is automatically satisfied by $N_{\mathrm{gen}} = \dim\mathbf{3}_f = 3$:
    \[
    3 \times (1, 1, 1) \;\equiv\; (0, 0, 0) \pmod 3.
    \]
\end{enumerate}
\end{theorem}

\begin{table}[htbp]
\centering
\caption{Exact 3-generation Standard Model chiral fermion and WIMP dark matter spectrum in $\R[\Gstar]$. Every SM multiplet has $\iota(\chi) = (-1) \times (\pm 1\text{ or }0) \times 0 = 0$, forming purely complex $3$-generation blocks without lepton collapse.}
\label{tab:spectrum}
\resizebox{\textwidth}{!}{%
\begin{tabular}{@{}lccccccc@{}}
\toprule
\textbf{Multiplet} & \textbf{$2I$ Irrep} & \textbf{$\Sigma(216\times 3)$ Irrep} & \textbf{$\Delta(27)_f$} & \textbf{$q_{\mathrm{PQ}}$} & \textbf{$\C$-Deg} & \textbf{Weyl States} & \textbf{Coupled Higgs Block} \\
\midrule
$Q_L = (u_L, d_L)$ & $2I[1]\;(2_{-1})$ & $\Sigma[18]\;(6_0, q_Y=0)$ & $\Delta[9]\;(3_0)$ & $0$ & $36$ & $18$ & --- \\
$u_R$ & $2I[2]\;(2_{-1}')$ & $\Sigma[11]\;(3_0, q_Y=+1)$ & $\Delta[9]\;(3_0)$ & $0$ & $18$ & $9$ & $H_u = (\Sigma[4],\mathrm{PQ}[0]) \in M_2(\C)_0$ \\
$d_R$ & $2I[2]\;(2_{-1}')$ & $\Sigma[10]\;(3_0, q_Y=-1)$ & $\Delta[9]\;(3_0)$ & $-1$ & $18$ & $9$ & $H_d = (\Sigma[3],\mathrm{PQ}[-1]) \in M_2(\C)_{\mp 1}$ \\
$L_L = (\nu_L, e_L)$ & $2I[1]\;(2_{-1})$ & $\Sigma[5]\;(2_{-1}, q_Y=0)$ & $\Delta[9]\;(3_0)$ & $+1$ & $12$ & $6$ & --- \\
$e_R$ & $2I[2]\;(2_{-1}')$ & $\Sigma[2]\;(1_0, q_Y=-1)$ & $\Delta[9]\;(3_0)$ & $0$ & $6$ & $3$ & $H_e = (\Sigma[3],\mathrm{PQ}[-1]) = H_d$ \\
$\nu_R$ & $2I[2]\;(2_{-1}')$ & $\Sigma[0]\;(1_{+1}, q_Y=0)$ & $\Delta[9]\;(3_0)$ & $+1$ & $6$ & $3$ & $H_\nu = (\Sigma[5],\mathrm{PQ}[0]) \in \quat_0$ \\
\midrule
\textbf{SM Total} & --- & Anomaly $= 0$ & $q_f = +1$ & $\equiv 0\bmod 3$ & $\mathbf{96}$ & $\mathbf{48}$ ($3\times 16$) & \textbf{3 Disjoint Color-Singlet Blocks} \\
\midrule
$\Phi_{\mathrm{DM}}$ (WIMP) & $2I[0]\;(1_{+1})$ & $\Sigma[6]\;(3_{+1} \in M_3(\R))$ & $\Delta[0]\;(1_{+1})$ & $0$ & $3$ & $3$ (Real Triplet) & Stable ($q_f = 0 \neq +1$) \\
\bottomrule
\end{tabular}%
}
\end{table}

\subsection{Structural Simplicity of \texorpdfstring{$\Gstar$}{G*}: The Clifford--Heisenberg Tower and \texorpdfstring{$E_8$}{E8} Golden--Eisenstein Duality}
\label{sec:e8_duality}

Beyond the bottom-up no-go theorems of Section~2.1, $\Gstar$ is not an arbitrary four-factor patchwork: it is completely determined by just two canonical algebraic objects corresponding to the two non-real Hurwitz division algebras ($\quat$ and $\C$) and the two arithmetic constructions of the $E_8$ root lattice $\Lambda_{E_8} \subset \R^8$:
\begin{enumerate}
    \item \textbf{The Internal Clifford--Heisenberg Normal Tower of $H = \Delta(27) \subset \SU(3)$:}
    Let $H \equiv \Delta(27) \cong H_3(\F_3) \subset \SU(3)$ denote the minimal faithful $3$D Weyl--Heisenberg group over $\F_3$. Its center is $Z(H) \cong \Z_3$ (order $3$), and its full normalizer inside $\SU(3)$---the single-qutrit Clifford group $\operatorname{Cl}_1(\F_3) = \{U \in \SU(3) : U H U^\dagger = H\}$, where the symplectic outer automorphism group $\operatorname{Sp}(2,\F_3) = \SL(2,\F_3) \cong 2T_W$ (order $24$) acts on the discrete phase space $H / Z(H) \cong \F_3^2$---is precisely the Hessian group:
    \begin{equation}
    N_{\SU(3)}(H) \;=\; H \rtimes \SL(2,\F_3) \;=\; \Delta(27)_c \rtimes 2T_W \;\cong\; \Sigma(216 \times 3) \qquad (\text{order } 27 \times 24 = 648).
    \end{equation}
    The three internal factors of $\Gstar$ are therefore the three canonical levels of the normal series of $H = \Delta(27)$ in $\SU(3)$,
    \begin{equation}
    \underbrace{Z(H)}_{\Z_3^{\mathrm{PQ}}\text{ (Peccei--Quinn grading)}} \;\triangleleft\; \underbrace{H}_{\Delta(27)_f\text{ (Flavor translations)}} \;\triangleleft\; \underbrace{N_{\SU(3)}(H)}_{\Sigma(216 \times 3)\text{ (Color + Weak gauge)}} \;\subset\; \SU(3),
    \end{equation}
    reducing the entire group $\Gstar$ to the maximal McKay-$E_8$ subgroup $2I \subset \SU(2)$ paired with the Clifford--Heisenberg tower of $H = \Delta(27) \subset \SU(3)$:
    \begin{equation}
    \label{eq:clifford_tower}
    \Gstar \;=\; \underbrace{2I}_{\text{Maximal in }\SU(2)} \;\times\; \underbrace{N_{\SU(3)}(H) \;\times\; H \;\times\; Z(H)}_{\text{Clifford--Heisenberg tower of }H = \Delta(27)\text{ in }\SU(3)}.
    \end{equation}

    \item \textbf{The Quaternionic / Golden Construction over $\Z[\phi]$ ($E_8 \to H_4 \oplus H_4' \longleftrightarrow G_{\mathrm{spin}} = 2I$):}
    Over the golden-ratio ring $\Z[\phi]$ ($\phi = \frac{1+\sqrt{5}}{2}$), $\Lambda_{E_8} \subset \R^8 \cong \quat^2$ is isomorphic to the \emph{icosian ring} $\mathcal{I} = \Z[\phi][2I]$ generated by the $120$ unit quaternions of the binary icosahedral group $2I \subset \SU(2)$~\cite{MoodyPatera1993}. Under the non-crystallographic folding $\R^8 \to \R^4_\parallel \oplus \R^4_\perp$, the $240$ roots of $E_8$ project into two golden-scaled copies of the $120$ vertices of the $4$D $600$-cell:
    \begin{equation}
    |\Phi(E_8)| \;=\; 240 \;=\; 120\;(H_4) \;+\; 120\;(\phi H_4') \;=\; |2I| + |2I'|.
    \end{equation}
    Crucially, the Galois involution $\sigma : \sqrt{5} \mapsto -\sqrt{5}$ ($\phi \mapsto 1-\phi = -1/\phi$) that exchanges the parallel and perpendicular $4$D subspaces $H_4 \leftrightarrow H_4'$ is identical to the outer automorphism of $2I$ that exchanges its two faithful $2$D quaternionic spinor representations in Table~\ref{tab:spectrum}:
    \begin{equation}
    \mathbf{2}_{\mathrm{spin},L} = 2I[1]\;\;(\chi_1(\mathcal{C}_2) = \phi \in H_4) \qquad \longleftrightarrow \qquad \mathbf{2}_{\mathrm{spin},R} = 2I[2]\;\;(\chi_2(\mathcal{C}_2) = -1/\phi \in H_4').
    \end{equation}
    Left--right fermion chirality in $\Gstar$ is thus directly realized as the $H_4 \leftrightarrow H_4'$ Galois conjugation of $E_8 \to H_4 \oplus H_4'$.

    \item \textbf{The Complex / Eisenstein Construction over $\Z[\omega]$ ($E_8 \to (E_6 \times \SU(3)_f)/\Z_3 \longleftrightarrow \Sigma(216 \times 3) \times \Delta(27)_f \times \Z_3^{\mathrm{PQ}}$):}
    Over the Eisenstein integers $\Z[\omega]$ ($\omega = e^{2\pi i/3}$), $\Lambda_{E_8} \subset \R^8 \cong \C^4$ is the rank-$4$ Hermitean lattice whose $240$ minimal vectors form the \emph{Witting polytope} $3\{3\}3\{3\}3\{3\}3$ in $\C^4$~\cite{ShephardTodd1954,Coxeter1974}. Its complex automorphism group is the Shephard--Todd reflection group $G_{32} \subset \U(4)$ of order
    \begin{equation}
    |G_{32}| \;=\; 155{,}520 \;=\; 240 \times 648 \;=\; |\Phi(E_8)| \times |\Sigma(216 \times 3)| \;=\; 2\,|G_2|,
    \end{equation}
    and its $3$D equatorial hyperplane section $\C^3 \subset \C^4$ is the $72$-vertex \emph{Hessian polyhedron} $3\{3\}3\{3\}3$ (the $72$ roots of $E_6 \subset E_8$), whose exact complex reflection symmetry group is $G_{25} \cong \Sigma(216 \times 3)$ (order $648$). Simultaneously, under the maximal Eisenstein subalgebra $E_8 \supset (E_6 \times \SU(3)_f)/\Z_3$, the $248$-dimensional adjoint of $E_8$ decomposes as $\mathbf{248} = (\mathbf{78},\mathbf{1}) \oplus (\mathbf{1},\mathbf{8}) \oplus (\mathbf{27},\mathbf{3}_f) \oplus (\overline{\mathbf{27}},\overline{\mathbf{3}}_f)$, partitioning the $240$ roots of $E_8$ as:
    \begin{equation}
    240 \;=\; \underbrace{72}_{E_6\text{ roots }\leftrightarrow\;\Sigma(216\times 3)} \;+\; \underbrace{6}_{\SU(3)_f\;\leftrightarrow\;\Delta(27)_f} \;+\; \underbrace{81}_{(\mathbf{27},\mathbf{3}_f)\;\leftrightarrow\;\Delta(27)_f \times \Z_3^{\mathrm{PQ}}} \;+\; \underbrace{81}_{(\overline{\mathbf{27}},\overline{\mathbf{3}}_f)}.
    \end{equation}
    Under $E_6 \supset \SO(10) \times \U(1)_{\mathrm{PQ}}$ (whose discrete $\Z_3$ center is $\Z_3^{\mathrm{PQ}}$), each fundamental $\mathbf{27}$ branches into $\mathbf{27} = \mathbf{16}_{+1} \oplus \mathbf{10}_{-2} \oplus \mathbf{1}_0$, so the $81 = |\Delta(27)_f \times \Z_3^{\mathrm{PQ}}|$ chiral states of $(\mathbf{27},\mathbf{3}_f)$ encompass precisely the $3 \times 16 = 48$ chiral Weyl fermions of Table~\ref{tab:spectrum} and the $3 \times 3 \times 2 = 18$ Weyl--Heisenberg Higgs doublets ($N_{H,\mathrm{high}} = 18$) active above $M_{\mathrm{fl}}$.

    \item \textbf{Bypassing the Distler--Garibaldi Continuum $E_8$ No-Go Theorem:}
    In continuum Lie theory, Distler and Garibaldi~\cite{DistlerGaribaldi2010} proved that one cannot embed both $\Spin(4)$ (from $E_8 \to H_4 \oplus H_4'$) and three chiral generations of $\SU(3)_c \times \SU(2)_{W,L} \times \U(1)_Y$ (from $E_8 \to (E_6 \times \SU(3)_f)/\Z_3$) inside a single real form of the Lie algebra $\mathfrak{e}_8$, because the $\Z[\phi]$ and $\Z[\omega]$ coordinate structures on $\Lambda_{E_8}$ are incompatible as simultaneous Lie subalgebras of $\mathfrak{e}_8$ and leave the $(\overline{\mathbf{27}},\overline{\mathbf{3}}_f)$ mirror fermions vector-like. In $(P, \Gstar, Z, \prec_Z)$, the two dual arithmetic symmetries of $\Lambda_{E_8}$ instead act as a \emph{direct product of finite groups},
    \begin{equation}
    \Gstar \;=\; \underbrace{2I}_{\Lambda_{E_8}\text{ over }\Z[\phi]\;(H_4 \oplus H_4')} \;\times\; \underbrace{\Sigma(216 \times 3) \times \Delta(27)_f \times \Z_3^{\mathrm{PQ}}}_{\Lambda_{E_8}\text{ over }\Z[\omega]\;(\text{Witting }G_{32} / E_6 \times \SU(3)_f)},
    \end{equation}
    pairing the chiral $H_4 \oplus H_4'$ spinors with the chiral $(\mathbf{27},\mathbf{3}_f)$ matter sector without continuum mirror doubling.
\end{enumerate}

%=============================================================================
\section{Dynamical Foundations: $d=4$ Selection, Continuous Symmetry Restoration, and Quantum Amplitudes}
%=============================================================================

\subsection{Dual Dynamical Selection of $d=4$ Spacetime: Percolation and Melting the Lattice Freezing Transition}

A central challenge for any theory based on finite subgroups $H \subset G_{\mathrm{Lie}}$ is the \emph{lattice freezing transition}: because a finite subgroup has a nonzero nearest-neighbor trace gap
\begin{equation}
\label{eq:trace_gap}
\Delta E_{\min}(\chi) \;=\; 1 - \max_{g \in H \setminus \{1\}} \frac{\operatorname{Re}\chi(g)}{d_\chi} \;>\; 0,
\end{equation}
a single-channel lattice action freezes into a discrete crystal above the critical inverse coupling $\beta_c(\chi) = 1 / \Delta E_{\min}(\chi)$. From the exact character tables of $2I$, $\Sigma(216 \times 3)$, and $\Delta(27)_f$, the minimum freezing thresholds across all sectors are $\beta_c(2I[1]) = 5.24$, $\beta_c(\Sigma[7]) = 2.32$, $\beta_c(\Sigma[5]) = 2.00$, $\beta_c(\Sigma[18]) = 1.50$, $\beta_c(\Sigma[19]) = 1.33$, and $\beta_c(\Delta_{27}[9]) = 1.00$.

At the poset unification cutoff $\Lambda$, the bare unified coupling $\alpha_U^{-1}(\Lambda) = 116/3 = 38.6667$ (Section~\ref{sec:couplings}) fixes the total $\SU(3)_c$ Wilson coupling $\beta_{\mathrm{tot}}(\Lambda) = 2N_c / g_U^2(\Lambda) = 18.462$. On a rigid periodic lattice ($1$ permutation channel), $\beta_{\mathrm{tot}} = 18.462 > \beta_c$ would freeze the group. On the $(d \to d)$ causal poset $(P, G, Z, \prec_Z)$ ($m = n = d$), however, Eq.~\eqref{eq:product_law} sums coherently over $d!$ permutation channels $\sigma \in S_d$, partitioning the effective coupling per channel as $\beta_{\mathrm{channel}}(d) = \beta_{\mathrm{tot}}(\Lambda) / d!$:

\begin{table}[htbp]
\centering
\caption{Dual dynamical selection of $d = 4$ spacetime coordination from causal percolation ($k_c(d) < 3$ for a minimal $\Z_3$ triplet fiber) and lattice freezing melting ($\beta_{\mathrm{channel}}(d) = \beta_{\mathrm{tot}}(\Lambda)/d! < \min_\chi \beta_c(\chi) = 1.00$).}
\label{tab:d4_selection}
\resizebox{\textwidth}{!}{%
\begin{tabular}{@{}cccccl@{}}
\toprule
\textbf{Degree $d$} & \textbf{Channels $d!$} & \textbf{Joint $p_{\mathrm{joint}}(d)$} & \textbf{Percolation $k_c(d)$} & \textbf{Channel $\beta_{\mathrm{channel}}(d)$} & \textbf{Dynamical Phase Status} \\
\midrule
$d = 1$ & $1$ & $4.763 \times 10^{-7}$ & $1449.0 > 3$ & $18.4615 > 1.00$ & Sub-critical \& Frozen Crystal \\
$d = 2$ & $2$ & $2.025 \times 10^{-6}$ & $26.50 > 3$ & $9.2308 > 1.00$ & Sub-critical \& Frozen Crystal \\
$d = 3$ & $6$ & $1.868 \times 10^{-5}$ & $6.15 > 3$ & $3.0769 > 1.00$ & Sub-critical \& Frozen Crystal \\
$\mathbf{d = 4}$ & $\mathbf{24}$ & $\mathbf{2.456 \times 10^{-4}}$ & $\mathbf{2.83 < 3}$ & $\mathbf{0.7692 < 1.00}$ & \textbf{Super-Critical \& Unfrozen Lie Phase!} \\
$d = 5$ & $120$ & $3.433 \times 10^{-3}$ & $1.76 < 3$ & $0.1538 < 1.00$ & Super-critical \& Unfrozen \\
\bottomrule
\end{tabular}%
}
\end{table}

\begin{theorem}[Dual Dynamical Selection of $d=4$ and Continuous Lie Restoration]
\label{thm:d4_melting}
Coordination degree $d = 4$ is the unique minimal vertex degree for which the causal poset simultaneously (i) crosses the percolation threshold $k_c(4) = p_{\mathrm{joint}}(4)^{-1/8} = 2.83 < 3$ for a minimal $\Z_3$ triplet fiber, and (ii) reduces the per-channel coupling $\beta_{\mathrm{channel}}(4) = 18.462 / 24 = 0.7692$ below $\min_\chi \beta_c(\chi) = 1.00$. By Burnside's Theorem on the simple Wedderburn blocks of $\R[\Gstar]$, the $S_4$-symmetrized vertex operators $\overline{U}_4 = \frac{1}{24}\sum_{\sigma \in S_4} \rho(u_{\sigma(1)}\cdots u_{\sigma(4)})$ densely generate the continuous Lie algebras $\mathfrak{spin}(4)$ and $\mathfrak{su}(3)_c \oplus \mathfrak{su}(2)_{W,L} \oplus \mathfrak{u}(1)_Y$ with vanishing coarse-grained trace gap ($\Delta E_{\mathrm{coarse}} \to 0$).
\end{theorem}

Furthermore, because $Z(2I) = \{\pm 1\}$ enforces parity/$\mathcal{CPT}$ inversion and $2I \subset \SU(2)_{\mathrm{spin}}$ is an exact unitary $5$-design (whose character moments $\frac{1}{|2I|}\sum_{g\in 2I} |\chi_1(g)|^{2t}$ equal the continuous $\SU(2)$ Haar Catalan numbers $C_t = \frac{1}{t+1}\binom{2t}{t} = (1, 2, 5, 14, 42)$ identically up to polynomial degree $2t = 10$), all linear ($n=1$, dimension-5) and anisotropic ($2t \le 10$) Lorentz-invariance-violating (LIV) operators vanish identically ($\xi_1 = 0$), while the quadratic ($n=2$) dispersion scale $\Lambda = 5.024 \times 10^{13}\text{ GeV}$ exceeds the $13\text{ TeV}$ photon bound from LHAASO GRB 221009A ($E_{\mathrm{QG},2} > 6.0 \times 10^{11}\text{ GeV}$) by a factor of $83.7$.

\subsection{Unitary Quantum Amplitudes $e^{-iHt/\hbar}$ and Non-SUSY Bose--Fermi Balance}

\begin{theorem}[Unitary Quantum Evolution and Exact Bose--Fermi Supertrace Balance]
\label{thm:quantum}
On the real group algebra $\R[\Gthree]$ ($\dim_{\R} = 2{,}099{,}520$), left translation $L_g : h \mapsto gh$ is orthogonal ($L_g^T = L_{g^{-1}}$). Consequently:
\begin{enumerate}
    \item \textbf{Unitary Quantum Phase Emergence:} Because causal links $x \prec_P y$ are directed, the infinitesimal generator of directed causal transport is the skew-symmetric operator $\mathcal{A}(g) = \frac{1}{2}(L_g - L_{g^{-1}})$ ($\mathcal{A}^T = -\mathcal{A}$). Decomposing $\R[\Gthree]$ by Frobenius--Schur indicator yields:
    \begin{equation}
    \begin{aligned}
    \dim\bigl(\iota = +1,\; M_n(\R)\bigr) &= 4{,}680 \;(0.22\%), \qquad
    \dim\bigl(\iota = -1,\; M_n(\quat)\bigr) = 4{,}680 \;(0.22\%), \\
    \dim\bigl(\iota = 0,\; M_n(\C)\bigr) &= 2{,}090{,}160 \;(99.55\%).
    \end{aligned}
    \end{equation}
    Across the $99.78\%$ of $\R[\Gthree]$ belonging to complex and quaternionic blocks, the canonical complex structure $J^2 = -I$ identifies $\mathcal{A}(g) \equiv -iH(g)/\hbar$ for a self-adjoint Hamiltonian $H(g) = H(g)^\dagger$, so Stone's Theorem yields unitary quantum time evolution $U(t) = \exp(-i\int H\,dt/\hbar)$.
    \item \textbf{Non-Supersymmetric Bose--Fermi Balance:} Because the bosonic ($z_{\mathrm{spin}} = +1$, $\iota = +1$) and fermionic ($z_{\mathrm{spin}} = -1$, $\iota = -1$) sectors of $2I$ both have $\sum d_\chi^2 = 60$, tensoring with $\Sigma(216 \times 3) \times \Delta(27)_f$ enforces exact equality between strictly real and quaternionic dimensions:
    \begin{equation}
    \STr_{\iota \neq 0}\bigl(\R[\Gthree]\bigr) \;=\; \dim\bigl(\R[\Gthree]_{\iota = +1}\bigr) - \dim\bigl(\R[\Gthree]_{\iota = -1}\bigr) \;=\; 4{,}680 - 4{,}680 \;=\; 0,
    \end{equation}
    achieving Bose--Fermi vacuum cancellation without unobserved spacetime superpartners.
\end{enumerate}
\end{theorem}

%=============================================================================
\section{Gauge Couplings, Gravity, and Plancherel Hierarchy Scales}
\label{sec:couplings}
%=============================================================================

\subsection{The Kernel Scale-Separation Theorem and Feynman--Plancherel Power-Counting}

In $\Gstar = \Gthree \times \Z_3^{\mathrm{PQ}}$, every unbroken gauge boson ($g, W^\pm, Z, \gamma$), the graviton ($2I$), the Weyl--Heisenberg flavor shifts ($W_{p,q}$), and the up-type Higgs $H_u = (\Sigma[4], \mathrm{PQ}[0])$ carry $q_{\mathrm{PQ}} = 0$, so $\Z_3^{\mathrm{PQ}} \subseteq \ker(\rho)$. Pure gauge, flavon-threshold, and gravitational links therefore act faithfully on the quotient group $G_{\mathrm{eff}} = \Gstar / \Z_3^{\mathrm{PQ}} \cong \Gthree$ ($|\Gthree| = 2{,}099{,}520$, $|\Gthree^{\mathrm{ab}}| = 27$). Conversely, processes probing all $\Z_3^{\mathrm{PQ}}$ sectors---specifically $\nu_R \sim (\Sigma[0], \mathrm{PQ}[+1])$ and closed vacuum loops---sample $|\Gstar| = 3|\Gthree| = 6{,}298{,}560$.

In the poset partition function $\mathcal{Z}[P, \Gstar]$, each vertex $v \in V(\Gamma_{\min})$ of the minimal connected diagram $\Gamma_{\min}(\mathcal{O})$ mediating an operator $\mathcal{O}$ inserts the Plancherel delta function $\delta_{G_v}(h_v) = \frac{1}{|G_v|}\sum_{\chi \in \Irr(G_v)} d_\chi \chi(h_v)$ of its faithful quotient group $G_v = \Gstar / \ker(\rho_v) \in \{\Gthree^{\mathrm{ab}}, \Gthree, \Gstar\}$, motivating the leading-order Feynman--Plancherel power-counting relations $\mu(\mathcal{O})/M_Z = d_{\chi_{\mathrm{ext}}}^{-1} \prod_{v \in V(\Gamma_{\min})} |G_v|^{\pm 1}$ from $1$-vertex, $2$-vertex, and $6$-link topologies:
\begin{align}
\frac{M_{\mathrm{DM}}}{M_Z} &= \frac{M_{A^0}}{M_Z} = |\Gthree^{\mathrm{ab}}|^1 = 27 &&\Longrightarrow M_{\mathrm{DM}} = M_{A^0} = 2.4621\text{ TeV}, \label{eq:scale_mdm} \\
\frac{M_{\mathrm{fl}}}{M_Z} &= |\Gthree| \cdot |\Gthree^{\mathrm{ab}}| = 5.6687 \times 10^7 &&\Longrightarrow M_{\mathrm{fl}} = 5.169 \times 10^9\text{ GeV}, \label{eq:scale_mfl} \\
\frac{\Lambda}{M_Z} &= \frac{|\Gthree|^2}{d_{\Sigma[19]}} = \frac{|\Gthree|^2}{8} = 5.5100 \times 10^{11} &&\Longrightarrow \Lambda = 5.024 \times 10^{13}\text{ GeV}, \label{eq:scale_lambda} \\
\frac{\Lambda_\nu}{M_Z} &= |\Gthree| \cdot |\Gstar| = 24\,\frac{\Lambda}{M_Z} = 1.3224 \times 10^{13} &&\Longrightarrow \Lambda_\nu = 1.206 \times 10^{15}\text{ GeV}, \label{eq:scale_lambdanu} \\
\frac{E_\Lambda}{M_Z} &= |\Gstar|^{-2} = \frac{1}{9|\Gthree|^2} = 2.5206 \times 10^{-14} &&\Longrightarrow E_\Lambda = 2.2985\text{ meV}, \label{eq:scale_elambda} \\
\alpha_G(m_p) &= \frac{1}{|Z(2I)|}|\Gthree|^{-6} = \frac{1}{2}(2{,}099{,}520)^{-6} &&\Longrightarrow \alpha_G(m_p) = 5.838 \times 10^{-39}. \label{eq:scale_alphag}
\end{align}
Specifically: (i) $(p,q)=(0,0)$ modes $\Phi_{\mathrm{DM}}, A^0$ ($[P,Q]=1$) correspond to a $1$-vertex abelian Kirchhoff insertion ($G_v = \Gthree^{\mathrm{ab}}$); (ii) $(p,q)\neq(0,0)$ flavons correspond to a $2$-vertex mixed propagator ($\Gthree \times \Gthree^{\mathrm{ab}}$); (iii) $\Lambda$ corresponds to a $2$-vertex $\Gthree$ correlator on the $d_{\Sigma[19]} = 8$ gluon octet; (iv) $\Lambda_\nu$ links $H_\nu$ ($G_{v_1}=\Gthree$) to $\nu_R$ ($G_{v_2}=\Gstar$); (v) $1$-vertex vacuum tadpoles vanish by Kirchhoff conservation and $\STr = 0$, leaving the $2$-vertex bubble $|\Gstar|^{-2}$ for $E_\Lambda$; and (vi) $\alpha_G(m_p)$ closes a $6$-link 3D spatial cube on $\Gthree / Z(2I)$.

\subsection{Three-Step Flavon-Higgs Gauge Unification ($N_H: 18 \to 2 \to 1$) and $\alpha_{\mathrm{em}}^{-1}(M_Z) = 127.941$}

At the cutoff $\Lambda$, the Spectral Action trace across the $48$ Weyl fermions gives $\Tr(T_{3L}^2) = 6$ and $\Tr(Y^2) = 10$, enforcing bare unification $g_3^2(\Lambda) = g_2^2(\Lambda) = g_1^2(\Lambda) = \frac{5}{3}g_Y^2(\Lambda)$ with $\sin^2\theta_W(\Lambda) = \frac{6}{6+10} = \frac{3}{8} = 0.3750$. Furthermore, the bare inverse unified coupling is fixed by the sum of the $N_{\mathrm{quark}} = 36$ colored Weyl quark states and the ratio of the real gluon block dimension $\dim M_8(\R) = 64$ to the electroweak subalgebra dimension $\dim\R[2T_W] = 24$:
\begin{equation}
\label{eq:bare_alpha_u}
\alpha_U^{-1}(\Lambda) = N_{\mathrm{quark}} + \frac{\dim M_8(\R)}{\dim\R[2T_W]} = 36 + \frac{8}{3} = \frac{116}{3} = 38.6667 \quad (\text{vs.\ RG } 38.6117,\; 0.14\%\text{ error}).
\end{equation}
Running in three group-determined stages---(i) from $\Lambda = M_Z |\Gthree|^2 / 8$ to $M_{\mathrm{fl}} = M_Z |\Gthree| |\Gthree^{\mathrm{ab}}|$ with all $N_{H,\mathrm{high}} = 18$ Weyl--Heisenberg Higgs doublets active, (ii) from $M_{\mathrm{fl}}$ to the 1-vertex Kirchhoff threshold $M_{A^0} = M_Z |\Gthree^{\mathrm{ab}}| = 27\,M_Z = 2.4621\text{ TeV}$ with the $N_{H,\mathrm{mid}} = 2$ flavor-diagonal zero-mode doublets $(H_u^{(0,0)}, H_d^{(0,0)})$ active, and (iii) from $27\,M_Z$ to $M_Z$ with only the $N_{H,\mathrm{low}} = 1$ light Standard Model Higgs doublet $h^0$ active---predicts with zero free parameters:
\begin{equation}
\begin{aligned}
\sin^2\theta_W(M_Z)_{\mathrm{pred}} &= 0.23111 \quad (\text{exp } 0.23121 \pm 0.00004,\; 0.045\%\text{ error},\; -2.6\sigma), \\
\alpha_s(M_Z)_{\mathrm{pred}} &= 0.1170 \quad (\text{exp } 0.1179 \pm 0.0009,\; 0.76\%\text{ error},\; -1.0\sigma), \\
\alpha_{\mathrm{em}}^{-1}(M_Z)_{\mathrm{pred}} &= 127.941 \quad (\text{exp } 127.951 \pm 0.009,\; 0.008\%\text{ error},\; -1.1\sigma;\; \alpha_{\mathrm{em}}^{-1}(0) = 137.0261\text{ with }\Delta_{\mathrm{QED}}).
\end{aligned}
\end{equation}

%=============================================================================
\section{Fermion and Boson Mass Ratios (Yukawa Couplings)}
%=============================================================================

\subsection{The $\Z_3$-Graded Georgi--Jarlskog Color Law and $b\text{--}\tau$ Determinant Unification}

Although $d_R \sim (\Sigma[10], \mathrm{PQ}[-1])$ and $e_R \sim (\Sigma[2], \mathrm{PQ}[0])$ couple to the same color-singlet Higgs block $H_d = H_e = (\Sigma[3], \mathrm{PQ}[-1])$, their $\Sigma(216 \times 3)$ bilinears have Clebsch--Gordan branching fractions into $\Sigma[3](2_0)$ of $2/2 = 1$ (leptons) versus $2 / (6 \times 3) = 1/9 = N_c^{-2}$ (quarks, single-color isometric amplitude $N_c^{-1} = 1/3$). Across the three $\Z_3$ shift grades $p \in \{0, +1, -1\} \pmod 3$ of $\Delta(27)_f$ (3rd, 2nd, and 1st generations), the bare quark-to-lepton Yukawa ratio obeys
\begin{equation}
\label{eq:gj_color_law}
\frac{y_d^{(p)}(\Lambda)}{y_e^{(p)}(\Lambda)} \;=\; N_c^{-p} \;=\; 3^{-p}, \qquad p \in \{0, +1, -1\} \pmod 3.
\end{equation}
Because $\prod_{p \in \Z_3} 3^{-p} = 3^0 \cdot 3^{-1} \cdot 3^{+1} = 1$, Eq.~\eqref{eq:gj_color_law} enforces exact determinant unification $\det M_d(\Lambda) = \det M_e(\Lambda)$. Using the quenched pure-gluon Yang--Mills anomalous dimension exponent $\gamma_m^{(0)} / (2 b_0^{\mathrm{YM}}) = 4 / (\frac{11}{3}N_c) = 4/11$ of the $\Sigma[19](\mathbf{8}_{+1}) \in M_8(\R)$ color-link dressing as an ansatz yields $R_{d/e}^{\mathrm{YM}}(M_Z) = (\alpha_s(M_Z)/\alpha_U(\Lambda))^{4/11} = 1.7313$ (vs.\ experimental $(m_d m_s m_b / m_e m_\mu m_\tau)_{M_Z}^{1/3} = 1.6709$, $+3.61\%$, whereas the unquenched $N_f = 6$ one-loop pure-gauge baseline prior to Yukawa damping is $R_{d/e}^{\mathrm{gauge}}(M_Z) = 2.170$) and gives for the three down-quark masses at $M_Z$: $m_b(M_Z) = 3.02\text{ GeV}$ (exp $2.85 \pm 0.03\text{ GeV}$, $+6.1\%$, $+5.8\sigma_{\mathrm{exp}}$), $m_s(M_Z) = 59.3\text{ MeV}$ (exp $53.5 \pm 4.6\text{ MeV}$, $+1.3\sigma$), and $m_d(M_Z) = 2.53\text{ MeV}$ (exp $2.67 \pm 0.19\text{ MeV}$, $-0.7\sigma$).

\subsection{Koide's Formula $Q_e = 2/3$ and the Clock Phase $\delta_e = 2/9\text{ rad}$ in $\R[\Delta(27)_f]$}

In $\R[\Delta(27)_f]$, the $p = 0$ Clock subalgebra decomposes into two simple Wedderburn blocks: $\R[\langle Q \rangle] \cong \R(\chi_{0,0}^{(f)}) \oplus \C(\chi_{0,1}^{(f)} \oplus \chi_{0,2}^{(f)})$. Equipartition of the Frobenius norm between the $\R$ and $\C$ blocks ($\|v_\R\|^2 = \|v_\C\|^2 \iff 2|v_1| = \sqrt{2}\,v_0$) in $\sqrt{Y_e} = v_0 I_3 + v_1 Q + \bar{v}_1 Q^2$ gives eigenvalues $\sqrt{m_k} = v_0\bigl(1 + \sqrt{2}\cos(\delta_e + 2\pi k/3)\bigr)$ and identically enforces Koide's formula:
\begin{equation}
\label{eq:koide}
Q_e \equiv \frac{m_e + m_\mu + m_\tau}{(\sqrt{m_e} + \sqrt{m_\mu} + \sqrt{m_\tau})^2} = \frac{2}{3} = 0.666667 \quad (\text{exp } 0.666661 \pm 0.000007,\; 0.0009\%\text{ error}).
\end{equation}
Fixing the Clock phase offset to the ratio of the spinor center order $|Z(2I)| = 2$ to the flavor abelianization order $|\Delta(27)_f^{\mathrm{ab}}| = 9$,
\begin{equation}
\label{eq:brannen_phase}
\delta_e \;=\; \frac{|Z(2I)|}{|\Delta(27)_f^{\mathrm{ab}}|} \;=\; \frac{2}{9}\text{ rad} \;=\; 0.222222\text{ rad} \qquad (\text{exp } 0.222230(8)\text{ rad},\; 0.003\%\text{ error},\; 0.9\sigma),
\end{equation}
predicts $m_e = 0.51096(4)_\tau\text{ MeV}$ ($0.007\%$ error, $1.0\sigma_\tau$) and $m_\mu = 105.652(7)_\tau\text{ MeV}$ ($0.006\%$ error, $0.9\sigma_\tau$) from $m_\tau = 1776.86(12)\text{ MeV}$.

\subsection{Up-Quark Hierarchy, $\tan\beta = |A_5| = 60$, Light Higgs, and Heavy 2HDM Quartet}

At $\mu = M_Z$, the running top-to-bottom mass ratio matches the order of the bosonic icosahedral rotation group $|2I / Z(2I)| = |A_5| = 60$ ($\tan\beta(M_Z)_{\mathrm{exp}} = m_t(M_Z)/m_b(M_Z) = 60.25 \pm 0.80$, $0.41\%$ error). Because $H_u = (\Sigma[4], \mathrm{PQ}[0])$ has $q_{\mathrm{PQ}} = 0$, the up-quark hierarchy scales quadratically with the Clock eigenvalues: $m_c(M_Z) = m_t(M_Z)(m_\mu/m_\tau)^2 = 0.594\text{ GeV}$ (exp $0.619 \pm 0.017\text{ GeV}$) and $m_u(M_Z) = \frac{1}{3}m_c(M_Z)(m_e/m_\mu) = 0.94\text{ MeV}$ (exp $1.27 \pm 0.19\text{ MeV}$).

In the gauge/Higgs boson sector, combining $\sin^2\theta_W(M_Z)_{\mathrm{pred}} = 0.23111$ with the 1-loop custodial top-quark shift $c_W^2\Delta\rho_t^{(1)} = c_W^2\frac{3 G_F m_t^2}{8\sqrt{2}\pi^2} = 0.00719$ yields $m_W = 80.332\text{ GeV}$ ($0.056\%$ error, $-3.7\sigma$ at pure 1-loop top level; $80.379\text{ GeV}$ when the full SM $\overline{\mathrm{MS}}\to\text{on-shell}$ shift $\Delta\hat{s}_W^2 = 0.00806$ is included, vs.\ exp $80.377 \pm 0.012\text{ GeV}$), while imposing the 2-loop phase-transition criticality boundary condition $\lambda(M_{\mathrm{fl}}) = 0$ at the flavon-Higgs condensation scale $M_{\mathrm{fl}} = M_Z |\Gthree| |\Gthree^{\mathrm{ab}}| = 5.169 \times 10^9\text{ GeV}$ predicts the light Standard Model Higgs boson mass:
\begin{equation}
m_{h^0}^{\mathrm{pred}}\bigl(\lambda(M_{\mathrm{fl}}) = 0\bigr) \;=\; 125.10\text{ GeV} \qquad (\text{exp } 125.25 \pm 0.17\text{ GeV},\; 0.12\%\text{ error},\; 0.9\sigma).
\end{equation}
Simultaneously, the orthogonal Type-II 2HDM combination $\Phi_\perp = -\sin\beta\,H_d^{(0,0)} + \cos\beta\,H_u^{(0,0)}$ acquires its mass at the 1-vertex abelian Kirchhoff threshold $M_{A^0} = M_Z |\Gthree^{\mathrm{ab}}| = 27\,M_Z = 2.4621\text{ TeV}$ (fixing the soft PQ-breaking scale $m_{12} = M_{A^0}/\sqrt{\tan\beta + \cot\beta} = 317.85\text{ GeV}$), predicting a nearly degenerate heavy Higgs quartet in the alignment limit ($\sin(\beta-\alpha) = 1 - \mathcal{O}(10^{-6})$):
\begin{equation}
\label{eq:heavy_2hdm}
\begin{aligned}
M_{A^0} &= 27\,M_Z = 2.4621\text{ TeV}, \qquad
M_{H^\pm} = \sqrt{M_{A^0}^2 + m_W^2} = 2.4634\text{ TeV}, \\
M_{H^0} &= \sqrt{M_{A^0}^2 + M_Z^2 \cos^2(2\beta)} = 2.4638\text{ TeV}.
\end{aligned}
\end{equation}
At $M_{A^0} = 2.4621\text{ TeV}$ and $\tan\beta = 60$, the 1-loop heavy Higgs $(A^0, H^0, H^\pm)$ penguin shifts the axial Wilson coefficient in $B_s^0 \to \mu^+\mu^-$ by $\Delta_P = -\Delta_S = -0.0206$, reducing $\operatorname{BR}(B_s^0 \to \mu^+\mu^-)$ by $-4.04\%$ from the SM value $3.66 \times 10^{-9}$ to $3.512 \times 10^{-9}$ (closer to the experimental world average $(3.34 \pm 0.27) \times 10^{-9}$, $0.64\sigma$), while the $\tan\beta$-independent charged-Higgs dipole shift $\Delta C_7(M_W) = -0.0153$ raises $\operatorname{BR}(B \to X_s \gamma)$ from $3.36 \times 10^{-4}$ to $3.583 \times 10^{-4}$ (vs.\ experimental $(3.49 \pm 0.19) \times 10^{-4}$, $0.49\sigma$), and tree-level flavon $\Delta S = 2$ kaon mixing $1/M_{\mathrm{fl}}^2 = 3.74 \times 10^{-20}\text{ GeV}^{-2}$ is $2.7 \times 10^9$ times below the UTfit limit.

%=============================================================================
\section{Dirac Neutrinos, WIMP Dark Matter, Strong CP, Baryogenesis, and Dark Energy}
%=============================================================================

\begin{enumerate}
    \item \textbf{Dirac Neutrino Spectrum, Complete PMNS Mixing Angles, and BBN $\Delta N_{\mathrm{eff}}$:} From $\Lambda_\nu = M_Z |\Gthree| |\Gstar| = |2T_W|\Lambda = 1.206 \times 10^{15}\text{ GeV}$, the atmospheric Dirac neutrino mass is $m_{\nu_3} = v_{\mathrm{ew}}^2 / \Lambda_\nu = 50.27\text{ meV}$ (exp $\sqrt{\Delta m_{31}^2} = 50.11 \pm 0.28\text{ meV}$, $0.33\%$ error). Shifting $\delta_e = 2/9$ by one elementary $2T_W$ discrete phase quantum $2\pi / |2T_W| = \pi/12$ ($\delta_\nu = 2/9 + \pi/12$) predicts $m_{\nu_2} = 8.84\text{ meV}$ and $m_{\nu_1} = 0.38\text{ meV}$, giving $\Delta m_{31}^2 = 2.527 \times 10^{-3}\text{ eV}^2$ ($0.6\sigma$), $\Delta m_{21}^2 = 7.79 \times 10^{-5}\text{ eV}^2$ ($1.8\sigma$), and $\sum m_{\nu_i} = 59.49\text{ meV}$ (Normal Ordering, strictly Dirac). Acting on the Walsh--Hadamard eigenbasis $V = F_3^\dagger$ ($\sin^2\theta_{23}^{(0)} = 1/2, \sin^2\theta_{12}^{(0)} = 1/3, \sin^2\theta_{13}^{(0)} = 0$), this same tetrahedral twist $\Delta\delta_{\ell\nu} = \pi/12$ predicts all three PMNS mixing angles and the Dirac CP phase without free parameters:
    \begin{equation}
    \label{eq:pmns_angles}
    \begin{aligned}
    \sin^2\theta_{13}^{\mathrm{PMNS}} &= \frac{1}{3}\sin^2\!\left(\frac{2\pi}{|2T_W|}\right) = \frac{2-\sqrt{3}}{12} = 0.02233 \quad (\text{exp } 0.02224 \pm 0.00057,\; 0.40\%), \\
    \sin^2\theta_{12}^{\mathrm{PMNS}} &= \frac{1}{3} - \sin^2\theta_{13}^{\mathrm{PMNS}} = \frac{2+\sqrt{3}}{12} = 0.3110 \quad (\text{exp } 0.307 \pm 0.012,\; 1.30\%), \\
    \sin^2\theta_{23}^{\mathrm{PMNS}} &= \frac{1}{2} + 2\sin^2\theta_{13}^{\mathrm{PMNS}} = \frac{5-\sqrt{3}}{6} = 0.5447 \quad (\text{exp } 0.546 \pm 0.021,\; 0.25\%), \\
    \delta_{\mathrm{CP}}^{\mathrm{PMNS}} &= \frac{3\pi}{2} - \left(\frac{2}{9} + \frac{\pi}{12}\right) = 242.27^\circ \quad (\text{exp } 232^\circ \pm 31^\circ,\; 0.33\sigma).
    \end{aligned}
    \end{equation}
    Furthermore, because $y_{\nu_3} = \sqrt{2}m_{\nu_3}/v_{\mathrm{ew}} = 2.89 \times 10^{-13}$, the Yukawa thermalization ratio at $T = M_Z$ is $\Gamma_{\mathrm{Yuk}}(\nu_R)/H(M_Z) = 2.59 \times 10^{-11} \ll 1$ (preventing Dirac sphaleron washout), and even UV-equilibrated $\nu_R$ decoupling at $M_{\mathrm{fl}}$ are diluted by $g_*(M_{\mathrm{fl}}) = 177.75$ to $\Delta N_{\mathrm{eff}}^{\mathrm{UV}}(\nu_R) = 3(10.75/177.75)^{4/3} = 0.0712 < 0.18$.
    \item \textbf{Automatically Stable $\Sigma[6](\mathbf{3}_{+1}) \in M_3(\R)$ WIMP Dark Matter:} In $\R[\Gstar]$, $\Sigma[6](\mathbf{3}_{+1})$ has $\iota_\Sigma = +1$. While tensoring with a spin-$1/2$ Weyl spinor $2I[1](\mathbf{2}_{-1})$ ($\iota_{2I} = -1$) gives $\iota_{\mathrm{tot}} = (-1)(+1)(+1) = -1 \implies M_3(\quat)$, the strictly real Wedderburn block $M_3(\R)$ ($\iota_{\mathrm{tot}} = (+1)(+1)(+1) = +1$) corresponds to the real spin-$0$ weak-triplet causal-link holonomy scalar $\Phi_{\mathrm{DM}} = (\Phi^+, \Phi^0, \Phi^-) \sim (2I[0](\mathbf{1}_{+1}), \Sigma[6](\mathbf{3}_{+1}), \Delta[0](\mathbf{1}_{+1}), \mathrm{PQ}[0]) \in M_3(\R)$ inside the $216$-dimensional bosonic subalgebra $\R[\Sigma(216)]$ (adding no extra unbalanced matter beta shift, $\Delta b_2 = 0$). Because $\Phi_{\mathrm{DM}}$ has central flavor triality $q_f = 0$ whereas all SM fermions have $q_f = +1$ (and all SM gauge/Higgs bosons are even in $Z(\Sigma(216))$), linear decays $\Phi_{\mathrm{DM}} \to \mathrm{SM}$ are forbidden to all orders, while 1-loop electroweak exchange lifts $\Phi^\pm$ above $\Phi^0$ by $\Delta M_\Phi = 172.17\text{ MeV} > m_{\pi^\pm}$ so $\Phi^\pm \to \Phi^0\pi^\pm$ decays in $\tau(\Phi^\pm) = 0.154\text{ ns} \ll t_{\mathrm{BBN}}$. At $M_{\mathrm{DM}} = M_Z |\Gthree^{\mathrm{ab}}| = 2.4621\text{ TeV}$, the parameter-free Gauge-Higgs commutator coupling $\frac{1}{2}g_2^2(H^\dagger H)\Phi^a\Phi^a$ in $\R[\Sigma(216)]$ ($\lambda_{H\Phi} = g_2^2 = 0.425$) with self-consistent freeze-out ($x_f = 25.51$) and non-perturbative Hulth\'en Sommerfeld enhancement ($\langle S_{\mathrm{eff}}\rangle = 1.227$) yields $\Omega_{\mathrm{DM}}^{\mathrm{Somm}} h^2 = 0.1190$ (Planck exp $0.1200 \pm 0.0012$, $0.82\%$ error; tree-level benchmark $0.1208$) and spin-independent nucleon cross section $\sigma_{\mathrm{SI}}(\Phi^0 N) = 3.66 \times 10^{-46}\text{ cm}^2$ ($2.0\times$ below the current LZ bound $7.40 \times 10^{-46}\text{ cm}^2$ at $2.46\text{ TeV}$, and $130\times$ above the neutrino floor).
    \item \textbf{Tree-Level Strong CP Vanishing and Jarlskog Invariant $J_{\mathrm{CP}}$:} Under the flavon inversion condition $c_{-p,-q} = c_{p,q}^*$ in the positive-definite Clock vacuum (Section~\ref{sec:critiques}, Critique~9), Hermitian Yukawa elements $M_q = M_q^\dagger \in \R[\Delta(27)_f]$ satisfy $\bar{\theta}_{\mathrm{QCD}}^{\mathrm{tree}} = \arg\det(M_u M_d) = 0$ at tree level while accommodating full off-diagonal CKM CP violation with $\delta_{\mathrm{CKM}} = \pi/3 + 1/9\text{ rad} = 66.37^\circ$ (exp $65.5^\circ \pm 1.5^\circ$) and $J_{\mathrm{CP}} = 2.93 \times 10^{-5}$ (exp $3.08 \pm 0.15 \times 10^{-5}$, $-1.0\sigma$).
    \item \textbf{Sphaleron Selection Rule and Baryon Asymmetry $\eta_B$:} Under $Z(\Delta(27)_f) \times \Z_3^{\mathrm{PQ}}$, $\Delta B = 1$ proton decay ($qqq\ell$, $q_f \equiv 1 \bmod 3$) and $\Delta B = 2$ dinucleon oscillation ($(u_R d_R d_R)^2$, $q_{\mathrm{PQ}} \equiv 2 \bmod 3$) are forbidden to all orders, whereas the $3$-generation electroweak sphaleron $\prod_{g=1}^3 (Q_L^{(g)})^3 L_L^{(g)}$ ($\Delta B = \Delta L = 3$) has $(q_f, q_{\mathrm{PQ}}) = (12, 3) \equiv (0, 0) \bmod 3$ and is identically allowed, predicting $\eta_B = (\Tr Y^2)_{\mathrm{tot}}^2 |\Gthree^{\mathrm{ab}}| / |\Gthree|^2 = 100 \times 27 / (2{,}099{,}520)^2 = 6.125 \times 10^{-10}$ (Planck exp $6.12 \pm 0.04 \times 10^{-10}$, $0.08\%$ error).
    \item \textbf{Cosmological Constant and Dark Energy Scale:} Bulk zero-point energy cancels at tree level by vertex Kirchhoff conservation (Theorem~\ref{thm:kirchhoff}), while the residual 2-vertex full-group vacuum loop gives $E_\Lambda = \rho_\Lambda^{1/4} = M_Z / |\Gstar|^2 = 2.2985\text{ meV}$ (exp $2.246 \pm 0.015\text{ meV}$, $2.34\%$ error) and $\rho_\Lambda / M_{\mathrm{Pl}}^4 = \frac{1}{4\cdot 3^8}(M_Z/m_p)^4 |\Gthree|^{-20} = 1.256 \times 10^{-123}$ (exp $1.145 \times 10^{-123}$).
\end{enumerate}

%=============================================================================
\section{Summary Table and Experimental Falsifiers}
%=============================================================================

\subsection{Six Sharp Experimental Falsifiers}

Because $\Gstar$ contains no continuous adjustable parameters, its discrete selection rules and threshold scales yield six sharp experimental falsifiers:
\begin{enumerate}
    \item \textbf{Zero Neutrinoless Double-Beta Decay ($0\nu\beta\beta = 0$):} Unbroken $Z(\Delta(27)_f) \cong \Z_3$ forbids $\Delta L = 2$ Majorana neutrino masses to all orders; any confirmed $0\nu\beta\beta$ signal (LEGEND-1000, nEXO, CUPID) falsifies the theory.
    \item \textbf{Normal Neutrino Ordering with $\sum_{i=1}^3 m_{\nu_i} = 59.5\text{ meV}$:} Inverted ordering is ruled out and $\sum m_\nu$ is pinned at $59.5\text{ meV}$ (Section~6), testable by DESI, Euclid, JUNO, DUNE, and CMB-S4.
    \item \textbf{Absolute Nucleon Stability ($p \to e^+\pi^0 = 0$, $n \leftrightarrow \bar{n} = 0$):} Protected to all orders by $Z(\Delta(27)_f) \times \Z_3^{\mathrm{PQ}}$, testable at Hyper-Kamiokande and DUNE.
    \item \textbf{No Light Peccei--Quinn Axion:} Because $\Z_3^{\mathrm{PQ}}$ is a discrete symmetry rather than a continuous Peccei--Quinn $\U(1)_{\mathrm{PQ}}$ with a pseudo-Goldstone boson, the framework predicts null results in light-axion haloscopes and helioscopes (ADMX, IAXO).
    \item \textbf{A $2.4621\text{ TeV}$ Real $\SU(2)_{W,L}$ Triplet WIMP ($\Phi_{\mathrm{DM}} \in M_3(\R)$):} Protected to all orders by $Z(\Delta(27)_f)$ with $\sigma_{\mathrm{SI}}(\Phi^0 N) = 3.66 \times 10^{-46}\text{ cm}^2$ (Section~6), directly testable by LZ/XLZD and via monochromatic $2.46\text{ TeV}$ $\gamma$-ray lines ($\Phi^0\Phi^0 \to \gamma\gamma, \gamma Z$) at CTA.
    \item \textbf{A Heavy Type-II 2HDM Quartet ($A^0, H^\pm, H^0$) at $2.4621\text{--}2.4638\text{ TeV}$ with $\tan\beta = 60$:} Sitting just above the current $2.0\text{--}2.1\text{ TeV}$ LHC Run-2/Run-3 exclusion boundary (Eq.~\eqref{eq:heavy_2hdm}), this quartet is directly testable via $\tan\beta$-enhanced $b\bar{b} \to H^0/A^0 \to \tau^+\tau^-, b\bar{b}$ and $gb \to tH^\pm \to tb, t\tau\nu$ production at the HL-LHC and FCC-hh.
\end{enumerate}

\newpage
\subsection{Summary of Parameter-Free Predictions}

Table~\ref{tab:master} compiles the 35 parameter-free structural and quantitative predictions derived in Sections~3--6 from the invariants of $\Gstar$ and compares them against current experimental values.

\begin{table}[!htbp]
\centering
\caption{Summary of 35 parameter-free predictions from $\Gstar$ compared with experiment.}
\label{tab:master}
\resizebox{\textwidth}{!}{%
\begin{tabular}{@{}llccc@{}}
\toprule
\textbf{Observable} & \textbf{Group-Algebraic Formula} & \textbf{Predicted} & \textbf{Experimental} & \textbf{Error} \\
\midrule
Spacetime Dimension $d$ & $k_c(d) < 3$ \& $\beta_{\mathrm{ch}}(d) = \beta_{\mathrm{tot}}/d! < 1$ & $d = 4$ & $d = 4$ & Exact \\
Bose--Fermi Supertrace & $\dim(\iota=+1) - \dim(\iota=-1)$ in $\R[\Gthree]$ & $4680 - 4680 = 0$ & $0$ (no SUSY) & Exact \\
Generations $N_{\mathrm{gen}}$ & $\dim\Delta[9] = 3$ \& $N_{\mathrm{gen}} \equiv 0 \bmod 3$ & $3$ & $3$ & Exact \\
Bare $\sin^2\theta_W(\Lambda)$ & $\Tr(T_{3L}^2) / (\Tr(T_{3L}^2) + \Tr(Y^2))$ & $3/8 = 0.3750$ & $0.3750$ (GUT) & Exact \\
Bare $\alpha_U^{-1}(\Lambda)$ & $N_{\mathrm{quark}} + \dim M_8(\R)/\dim\R[2T_W]$ & $116/3 = 38.667$ & $38.612$ (RG) & $0.14\%$ \\
Threshold $M_{\mathrm{fl}} / M_Z$ & $|\Gthree| \cdot |\Gthree^{\mathrm{ab}}|$ ($N_H: 2 \to 18$) & $5.669 \times 10^7$ & $5.891 \times 10^7$ & $3.77\%$ \\
Cutoff $\Lambda / M_Z$ & $|\Gthree|^2 / \dim\Sigma[19] = |\Gthree|^2 / 8$ & $5.510 \times 10^{11}$ & $5.563 \times 10^{11}$ & $0.95\%$ \\
$\sin^2\theta_W(M_Z)$ & 3-step RG ($N_H: 18 \to 2 \to 1$) & $0.23111$ & $0.23121(4)$ & $0.045\%$ ($-2.6\sigma$) \\
$\alpha_s(M_Z)$ & 3-step RG ($N_H: 18 \to 2 \to 1$) & $0.1170$ & $0.1179(9)$ & $0.76\%$ ($-1.0\sigma$) \\
$\alpha_{\mathrm{em}}^{-1}(0)$ & $\frac{8}{3}\alpha_U^{-1}(\Lambda) + \Delta_{\mathrm{RG}}^{(18\to 2\to 1)} + \Delta_{\mathrm{QED}}$ & $137.0261$ & $137.0360$ & $0.0072\%$ ($-1.1\sigma_{M_Z}$) \\
Gravity $\alpha_G(m_p)$ & $\frac{1}{|Z(2I)|}|\Gthree|^{-6}$ & $5.838 \times 10^{-39}$ & $5.906 \times 10^{-39}$ & $1.15\%$ \\
Planck Ratio $M_{\mathrm{Pl}}/m_p$ & $\sqrt{2}\,|\Gthree|^3$ & $1.309 \times 10^{19}$ & $1.301 \times 10^{19}$ & $0.58\%$ \\
$W$ Boson Mass $m_W$ & $M_Z \cos\theta_W^{\mathrm{OS}}$ (1-loop $c_W^2\Delta\rho_t^{(1)}$) & $80.332\text{ GeV}$ ($80.379\text{ SM}$) & $80.377(12)\text{ GeV}$ & $0.056\%$ ($-3.7\sigma$) \\
Light Higgs Mass $m_{h^0}$ & 2-loop criticality $\lambda(M_{\mathrm{fl}}) = 0$ & $125.10\text{ GeV}$ & $125.25(17)\text{ GeV}$ & $0.12\%$ ($-0.9\sigma$) \\
Heavy 2HDM $(A^0, H^\pm, H^0)$ & $M_{A^0} = M_Z |\Gthree^{\mathrm{ab}}| = 27\,M_Z$ ($\tan\beta = 60$) & $2.4621\text{--}2.4638\text{ TeV}$ & $> 2.0\text{ TeV}$ (LHC) & Testable \\
Rare $B$ ($\operatorname{BR}_{B_s\to\mu\mu}, \operatorname{BR}_{b\to s\gamma}$) & 1-loop $A^0/H^\pm$ penguins at $27\,M_Z, \tan\beta=60$ & $3.51\times 10^{-9}, 3.58\times 10^{-4}$ & $3.34(27)\times 10^{-9}, 3.49(19)\times 10^{-4}$ & $\le 0.6\sigma$ \\
Koide Invariant $Q_e$ & $\|v_\R\|^2 = \|v_\C\|^2$ in $\R[\langle Q \rangle]$ & $2/3 = 0.666667$ & $0.666661(7)$ & $0.0009\%$ ($+0.8\sigma$) \\
Lepton Phase $\delta_e$ & $|Z(2I)| / |\Delta(27)_f^{\mathrm{ab}}| = 2/9\text{ rad}$ & $0.222222\text{ rad}$ & $0.222230(8)\text{ rad}$ & $0.003\%$ ($-0.9\sigma$) \\
Masses $(m_e, m_\mu)$ & From $m_\tau$, $Q_e = 2/3$, $\delta_e = 2/9$ & $0.51096(4)_\tau, 105.652(7)_\tau\text{ MeV}$ & $0.51100, 105.658\text{ MeV}$ & $0.007\%$ ($\le 1.0\sigma_\tau$) \\
$(\det M_d / \det M_e)_{M_Z}^{1/3}$ & $R_{d/e}^{\mathrm{YM}}(M_Z) = (\alpha_s/\alpha_U)^{4/11}$ (quenched) & $1.7313$ ($2.170$ gauge) & $1.6709$ & $3.61\%$ \\
$(m_b, m_s, m_d)_{M_Z}$ & $3^{-p} R_{d/e}^{\mathrm{YM}}(M_Z) m_{e_p}(M_Z)$ (quenched $4/11$) & $3.02\text{ GeV}, 59.3, 2.53\text{ MeV}$ & $2.85\text{ GeV}, 53.5, 2.67\text{ MeV}$ & $+5.8\sigma, +1.3\sigma, -0.7\sigma$ \\
$\tan\beta(M_Z) = m_t/m_b$ & $|2I / Z(2I)| = |A_5|$ & $60.00$ & $60.25 \pm 0.80$ & $0.41\%$ ($-0.3\sigma$) \\
$(m_c, m_u)_{M_Z}$ & $m_t(m_\mu/m_\tau)^2$, $\frac{1}{3}m_c(m_e/m_\mu)$ & $0.594\text{ GeV}, 0.94\text{ MeV}$ & $0.619\text{ GeV}, 1.27\text{ MeV}$ & $-1.4\sigma, -1.7\sigma$ \\
$|V_{us}|, |V_{cb}|, |V_{ub}|$ & $\sqrt{m_d/m_s}$, $\frac{m_\mu/m_\tau}{\sqrt{2}}$, $(m_\mu/m_\tau)^2$ & $0.2234, 0.0416, 0.00354$ & $0.2250, 0.0418, 0.00373$ & $\le 1.4\sigma$ \\
$\delta_{\mathrm{CKM}}$ \& $J_{\mathrm{CP}}$ & $\pi/3 + 1/9\text{ rad}$ & $66.37^\circ, 2.93\times 10^{-5}$ & $65.5^\circ, 3.08\times 10^{-5}$ & $-1.0\sigma$ \\
Strong CP $\bar{\theta}_{\mathrm{QCD}}^{\mathrm{tree}}$ & Hermitian $\R[\Delta(27)_f]$ ($c_{-p,-q}=c_{p,q}^*$) & $0$ (tree ansatz) & $< 1.0 \times 10^{-10}$ & Tree-level ansatz \\
Neutrino Mass $m_{\nu_3}$ & $v_{\mathrm{ew}}^2 / (M_Z |\Gthree| |\Gstar|)$ & $50.27\text{ meV}$ & $50.11(28)\text{ meV}$ & $0.33\%$ ($+0.6\sigma$) \\
$\Delta m_{31}^2, \Delta m_{21}^2$ & $\delta_\nu = 2/9 + \pi/12\text{ rad}$ & $2.53\times 10^{-3}, 7.79\times 10^{-5}\text{ eV}^2$ & $2.51\times 10^{-3}, 7.41\times 10^{-5}\text{ eV}^2$ & $+0.6\sigma, +1.8\sigma$ \\
PMNS $(s_{13}^2, s_{12}^2, s_{23}^2)$ & $\frac{2-\sqrt{3}}{12}, \frac{2+\sqrt{3}}{12}, \frac{5-\sqrt{3}}{6}$ ($\Delta\delta_{\ell\nu}=\frac{\pi}{12}$) & $0.02233, 0.3110, 0.5447$ & $0.02224, 0.307, 0.546$ & $\le 0.3\sigma$ \\
PMNS Phase $\delta_{\mathrm{CP}}^{\mathrm{PMNS}}$ & $\frac{3\pi}{2} - (\frac{2}{9} + \pi/12)\text{ rad}$ & $242.27^\circ$ & $232^\circ \pm 31^\circ$ & $+0.33\sigma$ \\
WIMP DM $\Omega_{\mathrm{DM}} h^2$ \& $\sigma_{\mathrm{SI}}$ & $\Phi_{\mathrm{DM}} \in M_3(\R)$, $M_{\mathrm{DM}} = 27 M_Z$, $\lambda_{H\Phi}=g_2^2$ & $0.1190,\; 3.66\times 10^{-46}\text{ cm}^2$ & $0.1200(12),\; <7.4\times 10^{-46}\text{ cm}^2$ & $0.82\%$ ($-0.8\sigma$) \\
Baryon $\eta_B$ \& BBN $\Delta N_{\mathrm{eff}}$ & $(\Tr Y^2)_{\mathrm{tot}}^2 |\Gthree^{\mathrm{ab}}| / |\Gthree|^2$;\; $3(10.75/177.75)^{4/3}$ & $6.125 \times 10^{-10},\; \le 0.0712$ & $6.12(4) \times 10^{-10},\; < 0.18$ & $0.08\%$ ($+0.1\sigma$) \\
Dark Energy $E_\Lambda$ & $M_Z / |\Gstar|^2 = M_Z / (9|\Gthree|^2)$ & $2.2985\text{ meV}$ & $2.246(15)\text{ meV}$ & $2.34\%$ ($+3.5\sigma$) \\
\bottomrule
\end{tabular}%
}
\end{table}

\subsection{Structural Obstructions Avoided by the Framework}

We'd like to note several recurring structural obstructions in unification and quantum-gravity model building from which the present framework $(P, \Gstar, Z, \prec_Z)$ is free:
\begin{itemize}
    \item \textbf{Extra spatial dimensions, vacuum landscapes, and unstabilized moduli:} Spacetime coordination $d = 4$ is dynamically selected by causal percolation and lattice crystal melting (Theorem~\ref{thm:d4_melting}), leaving no compactification manifolds, continuous moduli fields, or landscape of competing vacua.
    \item \textbf{Unobserved spacetime supersymmetry:} Exact Bose--Fermi supertrace balance ($\dim(\iota=+1) = \dim(\iota=-1) = 4{,}680$, Theorem~\ref{thm:quantum}) is enforced algebraically by the binary icosahedral factor $2I$ without doubling the particle spectrum into superpartners.
    \item \textbf{Matterless geometry and unfixed quantization parameters:} Rather than quantizing pure geometry in isolation with an adjustable area-entropy parameter and appending gauge and matter fields externally, spacetime spinors, gauge bosons, Higgs doublets, and flavor generations emerge from a single finite group $\Gstar$ locked at every poset vertex by shared permutations $(\sigma, \tau) \in S_d^2$.
    \item \textbf{Pathological poset dominance and finite-subgroup lattice freezing:} Abelianization Kirchhoff conservation on $\Gstar^{\mathrm{ab}}$ (Theorem~\ref{thm:kirchhoff}) prevents high-degree combinatorial washout, while the $4! = 24$ vertex permutation channels reduce the per-channel coupling below the crystal-freezing threshold ($\beta_{\mathrm{channel}}(4) = 0.7692 < 1.00$).
    \item \textbf{Ad-hoc generation triplication and manual Yukawa couplings:} Both $N_{\mathrm{gen}} = 3$ and the CKM/PMNS mass and mixing hierarchies follow deductively from the Heisenberg flavor group $\Delta(27)_f$ and the discrete anomaly cancellation condition $3 \mid N_{\mathrm{gen}}$ (Theorem~\ref{thm:design}), rather than being postulated by hand or fitted via free continuous Yukawa parameters.
    \item \textbf{Leptoquark-mediated proton decay, doublet--triplet fine-tuning, and chiral mirror doubling:} Gauge unification arises from the Spectral Action trace on $\R[\Gstar]$ without leptoquark gauge bosons or colored Higgs triplets ($p \to e^+\pi^0 = 0$ to all orders under $Z(\Delta(27)_f)$), and the finite direct-product realization of the $\Z[\phi]$ and $\Z[\omega]$ symmetries of $\Lambda_{E_8}$ pairs chiral spinors with chiral matter without vector-like mirror doubling.
\end{itemize}

%=============================================================================
\section{Open Problems}
\label{sec:open_problems}
%=============================================================================

While the finite group $\Gstar$ is fixed within the ascending subgroup hierarchy by the representation-theoretic no-go theorems of Section~2 and reproduces the qualitative structure and leading-order magnitudes of Table~\ref{tab:master} without adjustable continuous parameters, four foundational and perturbative problems remain open:

\begin{enumerate}
    \item \textbf{Global Measure-Theoretic Continuum Limit as $|P| \to \infty$ and Emergent Boost Invariance vs.\ Finite Valence:}
    In pure Causal Set Theory (without the internal group fiber $\Gstar$), the Kleitman--Rothschild theorem~\cite{KleitmanRothschild1975} proves that as $|P| \to \infty$, almost all unconstrained finite posets collapse into pathological three-layer ``pancakes'' with macroscopic vertex coordination degree $d \sim |P|/4$. Within $(P, \Gstar, Z, \prec_Z)$, Theorem~\ref{thm:kirchhoff} and Theorem~\ref{thm:d4_melting} resolve the \emph{local} Kleitman--Rothschild pathology by suppressing high-degree hubs via the Kirchhoff selectivity floor $p_{\mathrm{joint}}(d) \le 1/|\Gstar^{\mathrm{ab}}| = 1/81$ while requiring $d \ge 4$ to percolate ($k_c(4) = 2.83 < 3$) and melt the lattice freezing transition ($\beta_{\mathrm{channel}}(4) = 0.7692 < 1.00$), pinning the microscopic vertex coordination degree at $d = 4$. Two non-perturbative continuum-limit problems remain: (i) proving that the path-integral sum over all locally finite posets $(P, \le_P)$ weighted by $\mathcal{Z}[P, \Gstar]$ dynamically self-organizes into a macroscopically 4-dimensional Alexandrov manifold (establishing that vertex coordination $d = 4$ induces macroscopic Hausdorff and spectral dimension $D = 4$ rather than a crumpled or branched-polymer 4-regular graph), and (ii) reconciling finite microscopic coordination $d = 4$ with the Bombelli--Henson--Sorkin theorem~\cite{BombelliHensonSorkin2009}---which proves that Poisson sprinklings into a continuum Minkowski manifold require infinite valence under exact boost invariance---by showing how the Galois parity involution $\sigma : \sqrt{5} \mapsto -\sqrt{5}$ on $2I \subset \Spin(4)$ and its unitary $5$-design isotropy (Section~\ref{sec:critiques}) restore continuous $\SL(2,\C)$ Lorentz boost invariance as an emergent infrared symmetry.

    \item \textbf{First-Principles Saddle-Point Derivation of the Plancherel Hierarchy Exponents and Continuum Gauge Dynamics:}
    In Section~\ref{sec:couplings} (Eqs.~\eqref{eq:scale_mdm}--\eqref{eq:scale_alphag}), the seven hierarchy scales ($M_{\mathrm{DM}}/M_Z = |\Gthree^{\mathrm{ab}}|$, $M_{\mathrm{fl}}/M_Z = |\Gthree|\cdot|\Gthree^{\mathrm{ab}}|$, $\Lambda/M_Z = |\Gthree|^2/8$, $\Lambda_\nu/M_Z = |\Gthree|\cdot|\Gstar|$, $E_\Lambda/M_Z = |\Gstar|^{-2}$, and $\alpha_G(m_p) = \frac{1}{2}|\Gthree|^{-6}$) and the baryon asymmetry formula $\eta_B = (\Tr Y^2)_{\mathrm{tot}}^2 |\Gthree^{\mathrm{ab}}|/|\Gthree|^2$ are motivated as the Plancherel weights of minimal connected $1$-vertex, $2$-vertex, and $6$-link topologies, and the coupling evolution from $\Lambda$ to $M_Z$ is evaluated using continuum Standard-Model renormalization-group equations. Deriving these dominant diagram weights uniquely from a master saddle-point theorem on $\mathcal{Z}[P, \Gstar]$, and demonstrating rigorously how massless continuum Yang--Mills gauge fields and continuous $\overline{\mathrm{MS}}$ $\beta$-function running emerge dynamically from the finite group algebra $\R[\Gstar]$, remain open mathematical-physics problems.

    \item \textbf{Higher-Loop, Threshold, and Yukawa Refinements of Leading-Order Quantitative Predictions:}
    Because all quantitative predictions in Table~\ref{tab:master} are evaluated at tree level or one loop without adjustable parameters or post-hoc resummation factors, several observables exhibit leading-order offsets that exceed experimental measurement uncertainties when theoretical truncation uncertainties are omitted:
    \begin{itemize}
        \item \emph{Electroweak, Neutrino, Cosmological, and Gravitational Offsets:} At one loop without threshold matching, $\sin^2\theta_W(M_Z) = 0.23111$ lies at $-2.6\sigma_{\mathrm{exp}}$ ($-0.045\%$), $\alpha_s(M_Z) = 0.1170$ at $-1.0\sigma_{\mathrm{exp}}$ ($-0.76\%$ in $N_f = 6$ prior to top-quark decoupling at $m_t$), and the pure 1-loop custodial $W$ mass $m_W = 80.332\text{ GeV}$ at $-3.7\sigma_{\mathrm{exp}}$ ($-0.056\%$ prior to SM bosonic/two-loop $\overline{\mathrm{MS}}\to\text{OS}$ corrections); the solar neutrino splitting $\Delta m_{21}^2 = 7.79 \times 10^{-5}\text{ eV}^2$ lies at $+1.8\sigma_{\mathrm{exp}}$ ($+5.2\%$); the tree-level dark energy scale $E_\Lambda = M_Z / |\Gstar|^2 = 2.2985\text{ meV}$ lies $+3.5\sigma_{\mathrm{exp}}$ ($+2.34\%$) above $2.246 \pm 0.015\text{ meV}$; and the tree-level gravitational coupling $\alpha_G(m_p) = \frac{1}{2}|\Gthree|^{-6} = 5.838 \times 10^{-39}$ ($M_{\mathrm{Pl}}/m_p = \sqrt{2}|\Gthree|^3 = 1.309 \times 10^{19}$) exhibits a $-1.15\%$ ($+0.58\%$) offset relative to $(m_p/M_{\mathrm{Pl}})^2 = 5.906 \times 10^{-39}$, which is large compared to the $2.2 \times 10^{-5}$ CODATA precision of $G_N$ and requires both a non-perturbative QCD derivation of $m_p/M_Z$ and a systematic $1/|\Gstar^{\mathrm{ab}}|$ loop expansion derived from $\mathcal{Z}[P, \Gstar]$.
        \item \emph{Quark Mass Running ($R_{d/e}(M_Z)$, $m_u(M_Z)$, $m_c(M_Z)$) and Radiative Strong CP Stability:} In Section~5.1, the $b\text{--}\tau$ determinant ratio $R_{d/e}^{\mathrm{YM}}(M_Z) = (\alpha_s(M_Z)/\alpha_U(\Lambda))^{4/11} = 1.7313$ ($+3.61\%$ above the experimental $1.6709$, giving $m_b(M_Z) = 3.02\text{ GeV}$ at $+5.8\sigma_{\mathrm{exp}}$) uses the quenched pure-gluon plaquette exponent $\gamma_m^{(0)}/(2b_0^{\mathrm{YM}}) = 4/11$ ($b_0^{\mathrm{YM}} = 11$), whereas the consistent unquenched $N_f = 6$ one-loop pure-gauge baseline ($b_3 = -7$) prior to Yukawa damping is $R_{d/e}^{\mathrm{gauge}}(M_Z) = 2.170$. Computing the coupled two-loop gauge and $\tan\beta = 60$ / $N_H = 18$ flavon-Higgs Yukawa running from $\Lambda$ to $M_Z$---both to determine whether Yukawa damping reduces the unquenched baseline $2.170$ to $1.6709$ and to refine $m_u(M_Z) = \frac{1}{3}m_c(M_Z)(m_e/m_\mu) = 0.94\text{ MeV}$ ($-1.7\sigma_{\mathrm{exp}}$) and $m_c(M_Z) = m_t(M_Z)(m_\mu/m_\tau)^2 = 0.594\text{ GeV}$ ($-1.4\sigma_{\mathrm{exp}}$)---as well as protecting the Hermitian Yukawa condition $M_q = M_q^\dagger$ against radiative loop shifts to $\bar{\theta}_{\mathrm{QCD}}$ (Critique~9), are left for future work.
    \end{itemize}

    \item \textbf{A Priori Formula Selection, Input Dependencies, and Look-Elsewhere Quantification:}
    Several entries in Table~\ref{tab:master} evaluate scaling laws using measured fermion mass ratios ($m_e, m_\mu, m_\tau, m_t$) or Standard-Model radiative shifts ($\Delta\alpha_{\mathrm{em}}(M_Z \to 0)$, $\Delta\rho_t$, two-loop Higgs $\beta_\lambda$) as inputs, or provide group-algebraic realizations of previously known empirical relations (Koide~\cite{Koide1983}, Brannen, and Gatto--Sartori--Tonin~\cite{GattoSartoriTonin1968}). Establishing a rigorous statistical significance that accounts for the look-elsewhere effect across the space of group invariants requires deriving every operator coefficient and threshold rule deductively from $\mathcal{Z}[P, \Gstar]$ without prior reference to experimental values.
\end{enumerate}

%=============================================================================
\section{Mathematical-Physical Discussion and Resolution of Foundational Objections}
\label{sec:critiques}
%=============================================================================

Replacing continuous spacetime manifolds and Lie groups with a discrete causal poset and a finite group raises legitimate mathematical and physical questions regarding continuum convergence, Lorentzian signature, decoupled representations, loop stability, cosmological relics, parameter counting, group-order minimality, and Strong CP. In this section, we examine nine foundational objections to the framework and detail how each is addressed within the algebraic structure of $\R[\Gstar]$.

\subsection{Critique 1: The Microscopic Poset Partition Function vs.\ Continuum Effective Action Bridge}

\paragraph{Objection.}
In Sections~1--5, the microscopic kinematics is discrete (a causal poset $(P, \prec)$ decorated by finite group elements $g_{xy} \in \Gstar$ obeying the ordered vertex product law), whereas the low-energy couplings and renormalization-group flows are evaluated using the continuum Connes--Chamseddine Spectral Action~\cite{ChamseddineConnes1997} and continuum $\overline{\mathrm{MS}}$ beta functions. Where is the explicit path-integral bridge showing that a purely combinatorial counting sum $\mathcal{Z}[P, \Gstar]$ converges to the continuum Einstein--Cartan and Yang--Mills actions without severe lattice-anisotropy artifacts?

\paragraph{Resolution (Plancherel Character Expansion and Unitary $t$-Design Theorem).}
On a causal poset $P$ whose vertices enforce the ordered product compatibility condition (Axiom~\ref{ax:conservation}), the microscopic weight of each elementary causal diamond (plaquette $\square$) with boundary holonomy $U_\square = g_1 g_2 g_3^{-1} g_4^{-1} \in \Gstar$ is the exact group delta function $\delta_{\Gstar}(U_\square)$. By the Plancherel theorem on the semisimple group algebra $\R[\Gstar]$, $\delta_{\Gstar}(U_\square)$ expands uniquely over the irreducible characters $\chi \in \operatorname{Irr}(\Gstar)$:
\begin{equation}
\label{eq:plancherel_partition}
\mathcal{Z}[P, \Gstar] = \sum_{\{g_{xy}\}} \prod_{\square \subset P} \delta_{\Gstar}(U_\square), \qquad
\delta_{\Gstar}(U_\square) = \frac{1}{|\Gstar|} \sum_{\chi \in \operatorname{Irr}(\Gstar)} d_\chi\,\chi(U_\square),
\end{equation}
which exponentiates via the heat-kernel character weights $\beta_\chi$ to $\exp\bigl(-\sum_\chi \beta_\chi(d_\chi - \operatorname{Re}\chi(U_\square)) + \cdots\bigr)$.
When the link holonomy $U_\square = \exp(i a^2 F_{\mu\nu})$ is expanded in powers of the fundamental cutoff spacing $a \sim \Lambda^{-1}$ up to dimension-4 continuum operators $\int d^4x\,\sqrt{g}\,\mathcal{L}_{\dim \le 4}$ (which correspond to polynomial moments of degree $2t \le 4$ in the group elements), a finite subgroup $H \subset \operatorname{U}(d)$ reproduces the exact continuum Haar measure of $\operatorname{U}(d)$ with zero anisotropy artifacts if and only if $H$ is a \emph{unitary $t$-design} for $t \ge 2$:
\begin{equation}
\frac{1}{|H|} \sum_{U \in H} |\Tr(U)|^{2t} \;=\; \int_{\operatorname{SU}(d)} |\Tr(U)|^{2t}\,d\mu_{\mathrm{Haar}}(U).
\end{equation}
Evaluating the character moments of the two non-abelian gauge/spin factors of $\Gstar$ directly from their character tables yields:
\begin{enumerate}
    \item \textbf{Spacetime Spin Factor $2I \subset \SU(2)_{\mathrm{spin}}$:} As noted in Section~3.1, the discrete moments $\frac{1}{120}\sum_{g \in 2I} |\chi_1(g)|^{2t}$ of $\chi_1 = \mathbf{2}_{\mathrm{spin},L}$ match the continuous $\SU(2)$ Catalan Haar moments $C_t$ identically for all $t \in \{1, \dots, 5\}$ ($2t \le 10$, numerical deviation $1.85 \times 10^{-13}$). Thus $2I$ is an \emph{exact unitary $5$-design}: rotational or Lorentz-violating lattice artifacts vanish not only at dimension $4$ ($2t = 4$), but up to dimension $10$, departing from continuous $\SU(2)$ first at $t = 6$ ($2t = 12$, where $131.6 \neq C_6 = 132$, suppressed by $(E/\Lambda)^8$).
    \item \textbf{Color/Weak Factor $\Sigma(216 \times 3) \subset \SU(3)_c$:} For the fundamental color triplet $\chi_7 = \mathbf{3}_c$, the discrete average $\frac{1}{648}\sum_{g \in \Sigma(216\times 3)} |\chi_7(g)|^{2t}$ equals the exact $\SU(3)$ Haar moments $(1, 2)$ for $t = 1, 2$ ($2t \le 4$, numerical deviation $1.47 \times 10^{-14}$). Thus $\Sigma(216 \times 3)$ is an \emph{exact unitary $2$-design} in $\SU(3)_c$, guaranteeing that the dimension-4 Yang--Mills action $\frac{1}{4g^2}\Tr(F_{\mu\nu}F^{\mu\nu})$ emerges from Eq.~\eqref{eq:plancherel_partition} with zero dimension-4 cubic or hypercubic anisotropy operators.
\end{enumerate}

\subsection{Critique 2: Compact \texorpdfstring{$2I \subset \Spin(4)$}{2I in Spin(4)} vs.\ Lorentzian \texorpdfstring{$\Spin(3,1) \cong \SL(2,\C)$}{Spin(3,1)} and the Bombelli--Henson--Sorkin Theorem}

\paragraph{Objection.}
The binary icosahedral group $2I$ is a finite, compact subgroup of Euclidean $\Spin(4) \cong \SU(2)_L \times \SU(2)_R$, whereas the non-compact Lorentzian spin group $\Spin(3,1) \cong \SL(2,\C)$ has no faithful finite-dimensional unitary representations. Furthermore, by the Bombelli--Henson--Sorkin theorem~\cite{BombelliHensonSorkin2009}, a discrete causal set generated by a Lorentz-invariant Poisson sprinkling into continuum Minkowski spacetime necessarily has infinite vertex valence. How can a poset with finite vertex coordination degree $d = 4$ and compact spin fiber $2I$ generate the non-compact boost commutator $[K_i, K_j] = -\epsilon_{ijk}J_k$ and avoid macroscopic Lorentz violation?

\paragraph{Resolution (Chiral Grading via $\sigma : \sqrt{5} \mapsto -\sqrt{5}$, Complexification, and Open Continuum Boost Problem).}
First, the Bombelli--Henson--Sorkin theorem~\cite{BombelliHensonSorkin2009} applies to a \emph{kinematic Poisson sprinkling of points into a pre-existing continuum Minkowski manifold}, where exact continuous boost invariance is imposed microscopically on a single static graph realization. In a non-perturbative path integral $\mathcal{Z}[P, \Gstar]$ over causal posets (as in 4D Causal Dynamical Triangulations and spin-foam quantum gravity), individual configurations have finite coordination degree $d = 4$, while whether continuous non-compact boost invariance and macroscopic Hausdorff/spectral dimension $D = 4$ are dynamically recovered in the coarse-grained superposition over poset histories remains an open problem (Section~\ref{sec:open_problems}, item~1). Second, because $2I \subset \Spin(4) \cong \SU(2)_L \times \SU(2)_R$ is finite and compact, no finite group or real group algebra $\R[2I]$ by itself generates non-compact Lorentz boosts (and the skew-symmetric transport operator $\mathcal{A}(g) = \frac{1}{2}(L_g - L_{g^{-1}})$ of Theorem~\ref{thm:quantum} satisfies $\mathcal{A}(g)^2 = \frac{1}{4}(L_{g^2} - 2I + L_{g^{-2}}) \neq -I$ in general, so it is not itself a complex structure). What $2I$ provides microscopically is the pair of distinct golden-ratio quaternionic spinor blocks $\quat(2I[1]) \oplus \quat(2I[2])$ ($\mathbf{2}_{\mathrm{spin},L} \oplus \mathbf{2}_{\mathrm{spin},R}$) interchanged by the \emph{Galois parity involution} $\sigma : \sqrt{5} \mapsto -\sqrt{5}$ of $\mathbb{Q}(\sqrt{5})$. Once these two real quaternion blocks are complexified ($\quat \otimes_\R \C \cong M_2(\C)$) and graded by $\sigma : 2I[1] \leftrightarrow 2I[2]$, they embed into the standard continuum Weyl representation $(1/2,0) \oplus (0,1/2)$ of $\mathfrak{sl}(2,\C) \cong \mathfrak{su}(2) \oplus i\,\mathfrak{su}(2)$:
\begin{enumerate}
    \item \textbf{Spatial rotation generators ($\sigma$-even):} Acting identically on $\mathbf{2}_L \oplus \mathbf{2}_R$ via anti-Hermitian generators $J_k = -\frac{i}{2}\operatorname{diag}(\sigma_k, \sigma_k)$ ($J_k^\dagger = -J_k$).
    \item \textbf{Boost generators ($\sigma$-odd):} Acting with opposite signs on $\mathbf{2}_L \oplus \mathbf{2}_R$ via the $\sigma$-grading $\operatorname{diag}(+I_2, -I_2)$ and multiplied by $i$, yielding Hermitian generators $K_k = \frac{1}{2}\operatorname{diag}(\sigma_k, -\sigma_k)$ ($K_k^\dagger = +K_k$).
\end{enumerate}
In this continuum Weyl embedding, the factor of $i$ ($i^2 = -1$) between anti-Hermitian $J_k$ and Hermitian $K_k$ produces the standard Lorentzian commutation relations:
\begin{equation}
\label{eq:lorentz_boost_algebra}
[J_i, J_j] \;=\; \epsilon_{ijk} J_k, \qquad
[J_i, K_j] \;=\; \epsilon_{ijk} K_k, \qquad
[K_i, K_j] \;=\; -\epsilon_{ijk} J_k,
\end{equation}
while the role specific to $2I$ is supplying the $\sigma : 2I[1] \leftrightarrow 2I[2]$ chirality grading and the exact unitary $5$-design isotropy (Section~3.1) of the compact rotation subgroup up to polynomial degree $2t = 10$.

\subsection{Critique 3: The Chiral Survival Theorem on \texorpdfstring{$\R[\Gstar]$}{R[G*]} (Why Exotic Irreps Decouple at \texorpdfstring{$\Lambda$}{Lambda})}

\paragraph{Objection.}
The group $\Gstar$ has $9 \times 24 \times 11 \times 3 = 7{,}128$ irreducible representations, including higher-spin fermionic irreps of $2I$ ($\chi_6$ of degree $4$, spin-$3/2$; $\chi_8$ of degree $6$, spin-$5/2$) and 9-dimensional irreps of $\Sigma(216 \times 3)$ ($\Sigma[22, 23]$). What dynamical principle prevents these non-Standard-Model irreps from appearing as light particles at $E \ll \Lambda$, rather than being discarded by hand?

\paragraph{Resolution (Georgi's Chiral Survival Theorem~\cite{Georgi1979} Applied to $\R[\Gstar]$).}
By the Survival Hypothesis of quantum field theory~\cite{Georgi1979}, any representation in $\R[\Gstar]$ that is vector-like (self-conjugate) under the unbroken symmetries above the electroweak scale—namely spatial parity $\sigma(2I)$, hypercharge $\Sigma(216 \times 3)^{\mathrm{ab}} \cong \Z_3$, and flavor triality $Z(\Delta(27)_f) \cong \Z_3$—admits a gauge-invariant bilinear mass term $M\bar{\psi}\psi$ in the Plancherel action~\eqref{eq:plancherel_partition} and acquires a cutoff mass $M \sim \Lambda = 5.02 \times 10^{13}\text{ GeV}$, decoupling completely from the low-energy spectrum. Only chirally protected representations remain massless down to $v_{\mathrm{ew}}$:
\begin{enumerate}
    \item \textbf{Elimination of Higher-Spin Rarita--Schwinger Fermions in $2I$:} Among the four fermionic irreps of $2I$ ($z_{\mathrm{spin}} = -1$: $\chi_1, \chi_2, \chi_6, \chi_8$), \emph{only} the spin-$1/2$ Weyl doublets $\chi_1(\mathbf{2}_L)$ and $\chi_2(\mathbf{2}_R)$ have irrational $\mathbb{Q}(\sqrt{5})$ characters ($\sigma(\chi_1) = \chi_2 \neq \chi_1$) and the minimal Plancherel action gap $\Delta E_{\min} = 1 - \cos(\pi/5) = 0.1910$. By contrast, the 4D spin-$3/2$ irrep $\chi_6 = \mathbf{2}_L \otimes \mathbf{2}_R$ ($d=4$) and the 6D spin-$5/2$ irrep $\chi_8$ ($d=6$) have \emph{strictly rational integer characters} across all 9 conjugacy classes of $2I$, so $\sigma(\chi_6) = \chi_6$ and $\sigma(\chi_8) = \chi_8$. Being parity self-conjugate (vector-like) with large Plancherel gaps $\Delta E_{\min} = 0.7500$ and $0.8333$, they acquire $\mathcal{O}(\Lambda)$ Dirac masses at the cutoff.
    \item \textbf{Elimination of the 9D Exotic Quarks $\Sigma[22, 23]$ in $\Sigma(216 \times 3)$:} Auditing the $\Z_3^{\mathrm{ab}}$ hypercharge orbits of all 24 irreps of $\Sigma(216 \times 3)$ reveals that $\Sigma[22](9_0)$ and $\Sigma[23](\bar{9}_0)$ have \emph{trivial $\Z_3^{\mathrm{ab}}$ orbits} $(22, 22, 22)$ and $(23, 23, 23)$: tensoring with the hypercharge generator $\lambda_Y \in \Sigma(216\times 3)^{\mathrm{ab}}$ leaves each invariant ($\chi_{22} \otimes \lambda_Y = \chi_{22}$), forcing their hypercharge to vanish ($Y = 0$). Because they are vector-like under $\U(1)_Y$ and pair with their complex conjugates in $\R[\Sigma(216 \times 3)]$, they form a vector-like Dirac pair at $M_9 \sim \Lambda$ and decouple.
    \item \textbf{Elimination of Abelian Flavor Singlets in $\Delta(27)_f$:} The 9 one-dimensional characters $\Delta[0..8]$ carry zero central flavor triality ($q_f = 0$, vector-like), leaving only the faithful 3D generation triplet $\mathbf{3}_f = \Delta[9]$ ($q_f = +1 \neq 0$).
\end{enumerate}
Consequently, out of all $7{,}128$ irreps of $\Gstar$, the $48 = 3 \times 16$ chiral Weyl fermions of the Standard Model (Table~\ref{tab:spectrum}) are the \emph{unique} chirally protected fermionic states in $\R[\Gstar]$.

\subsection{Critique 4: Naturalness of \texorpdfstring{$\delta_e = 2/9\text{ rad}$}{delta\_e = 2/9 rad} and Algebraic Homogeneity of \texorpdfstring{$\alpha_U^{-1}(\Lambda) = 36 + 8/3$}{alpha\_U\^-1(Lambda) = 36 + 8/3}}

\paragraph{Objection.}
Two algebraic expressions in Sections~4 and~5 invite skepticism regarding units and homogeneity: (i) in Eq.~\eqref{eq:brannen_phase}, why does a pure group-order ratio $\delta_e = |Z(2I)| / |\Delta(27)_f^{\mathrm{ab}}| = 2/9$ appear inside $\cos(\delta_e)$ in \emph{radians} rather than degrees or revolutions ($2\pi \times 2/9$)?; and (ii) in Eq.~\eqref{eq:bare_alpha_u}, why does $\alpha_U^{-1}(\Lambda) = 36 + 8/3 = 116/3$ add an integer fermion count ($36$) to a rational ratio of bosonic algebra dimensions ($64/24 = 8/3$)?

\paragraph{Resolution (Unitary Holonomy Exponential Map and Spectral-Action Trace Decomposition).}
Both expressions follow directly from standard functional analysis on $\R[\Gstar]$:
\begin{enumerate}
    \item \textbf{Why $\delta_e = 2/9$ is in radians:} In quantum mechanics, a hermitian generator $T_{\mathrm{Clock}}$ with unit-normalized spectrum ($\operatorname{spec}(T_{\mathrm{Clock}}) \in \{0, \pm 1\}$) generates a unitary holonomy operator via the operator exponential series
    \begin{equation}
    U_{\mathrm{Clock}} \;=\; \exp(i g_{\mathrm{Clock}} T_{\mathrm{Clock}}) \;=\; \sum_{n=0}^{\infty} \frac{(i g_{\mathrm{Clock}} T_{\mathrm{Clock}})^n}{n!}, \qquad
    \operatorname{Re}\langle U_{\mathrm{Clock}}\rangle \;=\; \cos(g_{\mathrm{Clock}}).
    \end{equation}
    Because the Taylor series $\cos x = \sum_{k=0}^\infty (-1)^k x^{2k}/(2k)!$ defines the cosine function with $x$ in \emph{radians} (the unique normalization in which $\frac{d}{dx}\sin x = \cos x$), any dimensionless coupling $g_{\mathrm{Clock}}$ multiplying a unit-normalized generator in $\exp(i g_{\mathrm{Clock}} T_{\mathrm{Clock}})$ automatically enters $\cos(g_{\mathrm{Clock}})$ in radians. Here $g_{\mathrm{Clock}} = |Z(2I)| / |\Delta(27)_f^{\mathrm{ab}}| = 2/9$ is the dimensionless coupling of the $\Z_2$ spinor center to each of the 9 abelian Clock/Shift character channels of $\Delta(27)_f^{\mathrm{ab}} \cong \Z_3 \times \Z_3$.
    \item \textbf{Why $\alpha_U^{-1}(\Lambda) = 36 + 8/3$ is homogeneous:} In the heat-kernel expansion of the Spectral Action $\Tr\,f(D_A/\Lambda)$ on the Hilbert space $\mathcal{H} = \mathcal{H}_F \oplus \mathcal{H}_B$~\cite{ChamseddineConnes1997}, the coefficient of the dimension-4 Yang--Mills curvature $\Tr(F_{\mu\nu}F^{\mu\nu})$ is the sum of two dimensionless traces: the fermionic trace over the colored Weyl spinor space $\mathcal{H}_F^{\mathrm{color}}$ ($\Tr_{\mathcal{H}_F}^{\mathrm{color}}(1) = N_{\mathrm{gen}} N_c N_{\mathrm{chirality}} N_{\mathrm{iso}} = 3 \times 3 \times 2 \times 2 = 36$) and the normalized bosonic bimodule trace of the gluon block $M_8(\R)$ relative to the weak isospin subalgebra $\R[2T_W]$ ($\dim_\R M_8(\R) / \dim_\R \R[2T_W] = 64/24 = 8/3$). Both terms are dimensionless spectral-action trace invariants on $\R[\Gstar]$.
\end{enumerate}

\subsection{Critique 5: All-Orders Factorization of Gauge/Yukawa-Weighted Bose--Fermi Supertraces}

\paragraph{Objection.}
Theorem~\ref{thm:quantum} establishes that the unweighted non-chiral Wedderburn dimension supertrace vanishes identically: $\operatorname{STr}_{\iota \neq 0}(1) = \sum_{\iota(\Pi)=\pm 1} \iota(\Pi)\dim_\R B_\Pi = 4680 - 4680 = 0$. In quantum field theory, however, 1-loop and higher-loop vacuum energy divergences $\Lambda^2 \operatorname{STr}(\mathcal{M}^2)$ and $\operatorname{STr}(\mathcal{M}^4 \ln\mathcal{M}^2)$ are weighted by mass matrices and gauge Casimir operators. Since bosons and fermions have different masses without supersymmetry, why don't gauge and Yukawa loop corrections reintroduce quadratic and quartic UV divergences?

\paragraph{Resolution (Exact Tensor-Product Factorization Across $\operatorname{Irr}(2I)$).}
Because the spacetime factor $2I$ commutes with the internal gauge and flavor factors $\Sigma(216 \times 3) \times \Delta(27)_f \times \Z_3^{\mathrm{PQ}}$, every Lorentz-invariant internal operator $\mathcal{O}_{\mathrm{int}} = 1_{2I} \otimes \mathcal{O}(\Sigma, \Delta)$—including arbitrary powers of the internal gauge Casimir operators $C_2(\Sigma)^k$, flavon mass matrices $\mathcal{M}_{\mathrm{fl}}^{2k}$, and cutoff masses $\mathcal{M}_\Lambda^{2k}$—acts as the identity $1_{2I}$ on the $2I$ factor. Since $\iota(\chi_{2I} \boxtimes \chi_\Sigma \boxtimes \chi_\Delta) = \iota(\chi_{2I})\iota(\chi_\Sigma)\iota(\chi_\Delta)$, the supertrace over the entire non-chiral sector ($\iota \neq 0$) factorizes identically:
\begin{equation}
\label{eq:factored_supertrace}
\operatorname{STr}_{\iota \neq 0}\!\bigl(1_{2I} \otimes \mathcal{O}(\Sigma, \Delta)\bigr) \;=\; \left(\sum_{\chi \in \operatorname{Irr}(2I)} \iota(\chi)\,d_\chi^2\right) \Tr_{\mathrm{int}}\!\bigl(\mathcal{O}(\Sigma, \Delta)\bigr) \;=\; (60 - 60)\,\Tr_{\mathrm{int}}(\mathcal{O}) \;\equiv\; 0.
\end{equation}
Equation~\eqref{eq:factored_supertrace} is a mathematical identity: \emph{any} operator $\mathcal{O}(\Sigma, \Delta)$ that depends only on internal gauge, Higgs, or flavor quantum numbers has a strictly vanishing supertrace at the cutoff $\Lambda$, because the bosonic ($A_5$, $\sum_{\iota=+1} d_\chi^2 = 1^2 + 3^2 + 3^2 + 4^2 + 5^2 = 60$) and fermionic ($2I \setminus A_5$, $\sum_{\iota=-1} d_\chi^2 = 2^2 + 2^2 + 4^2 + 6^2 = 60$) sectors of $2I$ share the exact same internal trace $\Tr_{\mathrm{int}}(\mathcal{O})$. A non-zero supertrace requires an operator that simultaneously breaks $2I$ spin grading and internal electroweak symmetry, which occurs only below the electroweak scale $v_{\mathrm{ew}}$ and is further protected at tree level by vertex Kirchhoff conservation (Theorem~\ref{thm:kirchhoff}).

\subsection{Critique 6: Cosmological \texorpdfstring{$\Z_3$}{Z\_3} Domain Walls and Pre-BBN Annihilation}

\paragraph{Objection.}
By the Zeldovich--Kobzarev--Okun theorem~\cite{Zeldovich1974}, the spontaneous breaking of discrete symmetries such as $\Z_3^{\mathrm{PQ}}$ and $\Delta(27)_f$ during cosmological phase transitions produces a network of topologically stable domain walls whose energy density scales as $\rho_{\mathrm{DW}} \propto a(t)^{-1}$, dominating over radiation ($\propto a^{-4}$) and matter ($\propto a^{-3}$) and destroying Big Bang Nucleosynthesis (BBN) and CMB isotropy. How does the theory avoid the cosmological domain-wall catastrophe?

\paragraph{Resolution (Explicit Vacuum Bias from the 3-Generation Anomaly and Soft $m_{12}$ Term).}
Cosmological domain walls are stable only if the discrete vacua remain exactly degenerate. In our framework, two built-in mechanisms explicitly lift the $\Z_3$ degeneracy and drive rapid pre-BBN domain-wall annihilation:
\begin{enumerate}
    \item Notice first that the center $Z(\Delta(27)_f) \cong \Z_3$ is \emph{never spontaneously broken} (all Higgs/flavon scalars reside in $\overline{\mathbf{3}}_f \otimes \mathbf{3}_f$ with $q_f = 0$), so $Z(\Delta(27)_f)$ forms zero domain walls.
    \item For the broken Peccei--Quinn $\Z_3^{\mathrm{PQ}}$ and abelian flavor $\Delta(27)_f^{\mathrm{ab}} \cong \Z_3 \times \Z_3$ symmetries, Theorem~\ref{thm:design} shows that the single-generation discrete QCD anomaly is non-zero ($\mathcal{A}_{\SU(3)_c^2\Z_3}^{(1)} = +1 \pmod 3$), and at the 1-vertex Kirchhoff threshold $M_{A^0} = 27\,M_Z = 2.4621\text{ TeV}$, the soft PQ-breaking mass $m_{12}^2 (H_u^\dagger H_d) + \mathrm{h.c.}$ with $m_{12} = 317.85\text{ GeV}$ explicitly biases the three $\Z_3^{\mathrm{PQ}}$ vacua by the energy density difference
    \begin{equation}
    \Delta V_{\mathrm{PQ}} \;\simeq\; \frac{m_{12}^2 v_{\mathrm{ew}}^2}{\tan\beta} \;=\; 1.02 \times 10^8\text{ GeV}^4.
    \end{equation}
\end{enumerate}
When the vacuum pressure $\Delta V_{\mathrm{PQ}}$ exceeds the surface tension force $\sigma_{\mathrm{DW}} / R_H(t)$, the domain walls collapse on a timescale $t_{\mathrm{coll}} \simeq \hbar\,\sigma_{\mathrm{DW}} / \Delta V_{\mathrm{PQ}}$:
\begin{equation}
t_{\mathrm{coll}}^{\mathrm{EW}} \;\simeq\; \frac{\hbar\,m_{12} v_{\mathrm{ew}}^2}{\Delta V_{\mathrm{PQ}}} \;=\; 1.24 \times 10^{-25}\text{ s}, \qquad
t_{\mathrm{coll}}^{\mathrm{fl}} \;\simeq\; \frac{\hbar\,M_{\mathrm{fl}} m_{12}^2}{\Delta V_{\mathrm{PQ}}} \;=\; 3.37 \times 10^{-18}\text{ s} \;\ll\; t_{\mathrm{BBN}} \sim 1\text{ s}.
\end{equation}
Even if one ignored $m_{12}$ entirely and retained only the non-perturbative QCD instanton bias $\Delta V_{\mathrm{QCD}} \simeq f_\pi^2 m_\pi^2 = 3.30 \times 10^{-4}\text{ GeV}^4$, the domain walls would still collapse at $t_{\mathrm{coll}}^{\mathrm{QCD}} = 3.84 \times 10^{-14}\text{ s} \ll 1\text{ s}$, leaving zero domain-wall relics at BBN.

\subsection{Critique 7: True Parameter Counting and Single-Scale Reduction to \texorpdfstring{$M_Z$}{M\_Z} Alone}

\paragraph{Objection.}
In Sections~4--7, while all dimensionless ratios and mixing angles are derived from $\Gstar$, absolute masses are quoted using $M_Z = 91.1876\text{ GeV}$ for gauge/Higgs/dark-matter/dark-energy scales, $v_{\mathrm{ew}} = 246.22\text{ GeV}$ for neutrinos, $m_t(M_Z) = 171.7\text{ GeV}$ for quarks, and $m_\tau = 1776.86\text{ MeV}$ for Koide's formula. Does the theory truly have only \emph{one} dimensionful scale (which merely defines the human unit $\text{GeV}$), or does it covertly rely on four independent mass anchors?

\paragraph{Partial Reduction and Current Limitations (Testing the $M_Z \to v_{\mathrm{ew}} \to m_t \to m_b \to m_\tau$ Chain).}
Every physical theory requires at least \emph{one} dimensionful reference value to define the arbitrary human unit of energy ($\text{GeV}$). Taking $M_Z = 91.1876\text{ GeV}$ as the dimensionful energy unit and imposing the discrete ans\"atze $y_t(M_Z) = 1$, $\tan\beta = |A_5| = 60$, and the quenched $4/11$ exponent $R_{d/e}^{\mathrm{YM}}(M_Z) = 1.7313$, one can test how closely the remaining anchors ($v_{\mathrm{ew}}, m_t, m_\tau$) are approximated:
\begin{enumerate}
    \item \textbf{Electroweak VEV $v_{\mathrm{ew}}$ from $M_Z$:} Using our one-loop predictions $\alpha_{\mathrm{em}}^{-1}(M_Z) = 127.941$ and $\sin^2\theta_W(M_Z) = 0.23111$ (Section~\ref{sec:couplings}), the electroweak gauge couplings at $M_Z$ are $g_2^2(M_Z) = 0.42525$ and $g_Y^2(M_Z) = 0.12781$, giving at tree level
    \begin{equation}
    v_{\mathrm{ew}}^{\mathrm{pred}}(M_Z) \;=\; \frac{2\,M_Z}{\sqrt{g_2^2(M_Z) + g_Y^2(M_Z)}} \;=\; 245.31\text{ GeV} \qquad (\text{exp } 246.22\text{ GeV},\; -0.37\%).
    \end{equation}
    \item \textbf{Top and Bottom Quark Masses $m_t(M_Z), m_b(M_Z)$ under the $y_t = 1, \tan\beta = 60$ Ans\"atze:} Assuming the unit-normalized top Yukawa coupling $y_t(M_Z) = 1$ on $H_u = (\Sigma[4], \mathrm{PQ}[0]) \in \quat$ and $\tan\beta = |A_5| = 60$ yields from $v_{\mathrm{ew}}^{\mathrm{pred}}$:
    \begin{equation}
    \begin{aligned}
    m_t^{\mathrm{pred}}(M_Z) &= \frac{v_{\mathrm{ew}}^{\mathrm{pred}}}{\sqrt{2}} = 173.46\text{ GeV} \quad (\text{exp } 171.7 \pm 1.5\text{ GeV},\; +1.02\%), \\
    m_b^{\mathrm{pred}}(M_Z) &= \frac{m_t^{\mathrm{pred}}(M_Z)}{|A_5|} = 2.891\text{ GeV} \quad (\text{exp } 2.85 \pm 0.03\text{ GeV},\; +1.44\%).
    \end{aligned}
    \end{equation}
    \item \textbf{Tau Lepton Mass $m_\tau$ from $M_Z$:} From $p = 0$ third-generation $b\text{--}\tau$ unification $y_b(\Lambda) = y_\tau(\Lambda)$ (Eq.~\eqref{eq:gj_color_law}) and the quenched pure-gluon ansatz $R_{d/e}^{\mathrm{YM}}(M_Z) = (\alpha_s(M_Z)/\alpha_U(\Lambda))^{4/11} = 1.7313$, the running and pole tau masses evaluate to:
    \begin{equation}
    \begin{aligned}
    m_\tau^{\mathrm{pred}}(M_Z) &= \frac{m_b^{\mathrm{pred}}(M_Z)}{R_{d/e}^{\mathrm{YM}}(M_Z)} = 1.6698\text{ GeV} \quad (\text{exp } 1.7462\text{ GeV},\; -4.37\%), \\
    m_\tau^{\mathrm{pole,pred}} &= m_\tau^{\mathrm{pred}}(M_Z)\bigl(1 + \delta_{\mathrm{QED}}^{\mathrm{pole}}\bigr) = 1699.18\text{ MeV} \quad (\text{exp } 1776.86\text{ MeV},\; -4.37\%).
    \end{aligned}
    \end{equation}
\end{enumerate}
Via Koide's formula ($Q_e = 2/3, \delta_e = 2/9\text{ rad}$), $m_\tau^{\mathrm{pole,pred}} = 1699.18\text{ MeV}$ yields $m_\mu^{\mathrm{pred}} = 101.03\text{ MeV}$ and $m_e^{\mathrm{pred}} = 0.4886\text{ MeV}$ ($-4.37\%$). Note, however, that a complete reduction to $M_Z$ alone is \emph{not yet achieved}: in the non-supersymmetric Type-II 2HDM one-loop renormalization-group equation above $M_{A^0}$,
\begin{equation}
16\pi^2 \frac{d\ln(y_b/y_\tau)}{d\ln\mu} \;=\; -8 g_3^2 + 2 g_1^2 + \frac{1}{2}y_t^2 + \frac{3}{2}y_b^2 - \frac{3}{2}y_\tau^2,
\end{equation}
the consistent unquenched $N_f = 6$ pure-gauge baseline ($-8g_3^2 + 2g_1^2$) prior to Yukawa damping is $R_{d/e}^{\mathrm{gauge}}(M_Z) = 2.170$ (which would give $m_\tau^{\mathrm{pole}} = 1355.8\text{ MeV}$, $-23.7\%$). One-loop Type-II 2HDM Yukawa running between $27\,M_Z$ and $M_{\mathrm{fl}}$ damps $2.170$ by only $\sim 5\%$ (to $\sim 2.06$), so reaching $1.7313$ or the experimental $1.6709$ requires an additional $\sim 20\%$ damping from the uncomputed $N_H = 18$ flavon-Higgs Yukawa couplings above $M_{\mathrm{fl}}$ (Section~\ref{sec:open_problems}, item~3). Finally, when assessing the pulls in Table~\ref{tab:master}, two symmetrical caveats apply: (i) pairs of rows such as $(\alpha_G, M_{\mathrm{Pl}}/m_p)$, $(E_\Lambda, \rho_\Lambda/M_{\mathrm{Pl}}^4)$, and $((Q_e,\delta_e), (m_e,m_\mu))$ are algebraically equivalent rather than independent tests; and (ii) a leading-order theoretical truncation uncertainty $\sigma_{\mathrm{th}} \sim \mathcal{O}(1/|\Gstar^{\mathrm{ab}}|) \sim 1.2\%$ cuts both ways---it explains why a $1.15\%$ tree-level offset in $\alpha_G(m_p)$ is natural at leading order despite the $2.2 \times 10^{-5}$ experimental precision of $G_N$, but it simultaneously implies that sub-per-mille one-loop agreements (such as $\sin^2\theta_W(M_Z)$ to $0.045\%$ or $\alpha_{\mathrm{em}}^{-1}(M_Z)$ to $0.008\%$) lie inside the one-loop theoretical uncertainty band until two-loop and threshold corrections are computed.

\subsection{Critique 8: Group-Order Minimality of \texorpdfstring{$\Gstar$}{G*} and Why \texorpdfstring{$\Delta(27)$}{Delta(27)} Rather Than \texorpdfstring{$T_7$}{T\_7} or \texorpdfstring{$\Z_9 \rtimes \Z_3$}{Z\_9 x| Z\_3}}

\paragraph{Objection.}
Section~2 analyzes the ascending sequence of four candidate groups $G_1 \subset G_2 \subset \Gthree \subset \Gstar$. Does the 5-Point Design Theorem (Theorem~\ref{thm:design}) merely select $\Gstar$ from a hand-picked list? In particular: (i) the Frobenius group $T_7 = \Z_7 \rtimes \Z_3$ of order $21 = 3 \cdot 1^2 + 2 \cdot 3^2 < 27$ also possesses a faithful irreducible complex $3$-dimensional representation ($\iota = 0$), and (ii) at order $27$, the other extraspecial group $M_{27} = \Z_9 \rtimes \Z_3$ (exponent $9$) also has a $\Z_3$ center, $|M_{27}^{\mathrm{ab}}| = 9$, and a faithful complex $3$D irrep spanning $M_3(\C)$. Why is $\Delta(27) = (\Z_3 \times \Z_3) \rtimes \Z_3$ (exponent $3$) selected?

\paragraph{Resolution (Schur's Center-Index Theorem, $\SL(2,\F_3)$ Automorphism Action, and Exponent-$3$ Symmetry).}
Each of the four factors of $\Gstar$ has the \emph{strictly minimal order over all finite groups} compatible with the requirements of Theorem~\ref{thm:design}, with the distinction between $T_7$, $M_{27} = \Z_9 \rtimes \Z_3$, and $\Delta(27) = (\Z_3 \times \Z_3) \rtimes \Z_3$ governed by three properties:
\begin{enumerate}
    \item \textbf{Schur's Center-Index Theorem ($d_\chi \mid [H : Z(H)]$) Eliminates Centerless $T_7$ (Order $21$):} By Schur's theorem, the degree $d_\chi$ of any irreducible representation of a finite group $H$ divides the index of its center: $d_\chi \mid |H / Z(H)|$. The Frobenius group $T_7$ achieves order $21$ only because its center is trivial ($|Z(T_7)| = 1$) and $|T_7^{\mathrm{ab}}| = 3$ (so $\overline{\mathbf{3}} \otimes \mathbf{3} = \mathbf{1}_0 \oplus \mathbf{1}_1 \oplus \mathbf{1}_2 \oplus \mathbf{3} \oplus \overline{\mathbf{3}}$ fails to decompose into $9$ one-dimensional characters). Imposing the physical design requirement of a non-trivial $\Z_3$ center $Z(H) \cong \Z_3$ (for $\SU(3)_c$ color $N$-ality confinement and for central flavor triality $Z(G_{\mathrm{fl}}) \cong \Z_3$ stabilizing the proton and WIMP) forces $3 \mid |Z(H)|$ and, by Schur's theorem $3 \mid (|H|/|Z(H)|)$, $9 \mid |H|$; since no group of order $9$ or $18$ has a $3$D irrep, $|H| = 27$ is the strict minimal order over all finite groups.
    \item \textbf{Why Color Requires $\Delta(27)_c$ (Exponent $3$) Rather Than $M_{27} = \Z_9 \rtimes \Z_3$ (Exponent $9$):} Up to isomorphism, there are two non-abelian groups of order $27$: $\Delta(27) \cong (\Z_3 \times \Z_3) \rtimes \Z_3$ (exponent $3$) and $M_{27} \cong \Z_9 \rtimes \Z_3$ (exponent $9$). Both have $|Z(H)| = 3$, $|H^{\mathrm{ab}}| = 9$, and faithful $3$D irreps spanning $M_3(\C)$. In the color sector, however, fusing the $8$ non-trivial characters of $H^{\mathrm{ab}}$ into a single irreducible $8$D real gluon octet $\mathbf{8}_{+1} \in M_8(\R)$ requires the weak group $2T_W \cong \SL(2,\F_3)$ (order $24$) to act transitively on the $8$ non-trivial central cosets of $H/Z(H) \cong \Z_3 \times \Z_3$. For $\Delta(27)_c$, every non-identity element has order $3$ and $\operatorname{Out}_Z(\Delta(27)_c) \cong \GL(2,\F_3)$ (order $48$) contains $2T_W \cong \SL(2,\F_3)$ acting transitively on the $8$ non-trivial cosets, forming $\Sigma(216 \times 3) \cong \Delta(27)_c \rtimes 2T_W$ (order $648$). By contrast, in $M_{27} = \Z_9 \rtimes \Z_3$, the elements of order dividing $3$ form a characteristic subgroup $\Omega_1(M_{27}) \cong \Z_3 \times \Z_3$ of order $9$: $6$ of the $8$ non-trivial cosets of $M_{27}/Z(M_{27})$ consist of elements of order $9$ while only $2$ consist of elements of order $3$, so $\operatorname{Aut}(M_{27})$ cannot mix them or host $2T_W \cong \SL(2,\F_3)$. In the flavor sector, on the other hand, $M_{27} = \Z_9 \rtimes \Z_3$ satisfies Weyl--Heisenberg completeness ($\dim\operatorname{span}\rho(M_{27}) = 9$) equally well, so selecting $\Delta(27)_f$ (exponent $3$, where Clock $Q$ and Shift $P$ enter symmetrically with $P^3 = Q^3 = I_3$) over $M_{27}$ (exponent $9$) is an additional symmetry choice between the two extraspecial groups of order $27$.
    \item \textbf{Minimal Chiral Spinor $5$-Design ($2I$, order $120$):} By the McKay ADE classification of finite subgroups of $\SU(2)$ ($C_n, 2D_{2n}, 2T, 2O, 2I$), the binary icosahedral group $2I$ of order $120$ ($E_8$) is the unique subgroup possessing two distinct faithful $2$D quaternionic spinor doublets ($\mathbf{2}_L \oplus \mathbf{2}_R$) interchanged by a quadratic Galois involution ($\Q(\sqrt{5})$) and forming an exact unitary $5$-design.
\end{enumerate}
The ascending chain $G_1 \subset G_2 \subset \Gthree \subset \Gstar$ then tests whether fewer than four minimal factors can share roles, which No-Go Theorems~\ref{thm:nogo1}--\ref{thm:nogo3} rule out by exact algebraic identities.

\subsection{Critique 9: What Is Assumed vs.\ Proved in the \texorpdfstring{$\R[\Delta(27)_f]$}{R[Delta(27)\_f]} Strong CP Mechanism?}

\paragraph{Objection.}
In Section~6 (item~3), tree-level Strong CP conservation $\bar{\theta}_{\mathrm{QCD}}^{\mathrm{tree}} = \arg\det(M_u M_d) = 0$ is verified on positive-definite Hermitian Yukawa elements $M_q = M_q^\dagger \in \R[\Delta(27)_f]$. However: (i) any positive-definite Hermitian matrix has $\arg\det = 0$ by linear algebra; (ii) in $\Gstar$, $Q_L \sim \Sigma[18]$ and $q_R \sim \Sigma[10,11]$ sit in different gauge irreps without a left--right gauge parity; and (iii) what protects $M_q = M_q^\dagger$ and $\bar{\theta}_{\mathrm{QCD}}$ against radiative corrections from gauge, 2HDM, and flavon loops?

\paragraph{Resolution (Flavon Inversion Ansatz, Weyl--Heisenberg Completeness, and Open Radiative Protection).}
It is important to separate what is an algebraic theorem of $\Delta(27)_f$, what is a tree-level symmetry ansatz on the flavon VEVs, and what remains an open radiative problem:
\begin{enumerate}
    \item \textbf{Tree-Level Flavon Inversion Ansatz ($c_{-p,-q} = c_{p,q}^*$) and Positive Clock Vacuum:} While $Q_L \sim (\Sigma[18], \Delta[9])$ and $q_R \sim (\Sigma[10,11], \Delta[9])$ carry distinct $\Sigma(216 \times 3)$ gauge irreps, both transform in the \emph{same} 3D flavor irrep $\mathbf{3}_f = \Delta[9]$ of $\Delta(27)_f$. Once the electroweak Higgs VEV contracts the gauge indices, the $3 \times 3$ flavor mass matrix $M_q = \sum_{p,q \in \Z_3^2} c_{p,q}^{(q)} W_{p,q}$ is Hermitian ($M_q^\dagger = M_q$) if the tree-level flavon VEVs obey the inversion condition $c_{-p,-q}^{(q)} = (c_{p,q}^{(q)})^*$ under $W_{p,q}^\dagger = W_{-p,-q}$, and $\arg\det(M_u M_d) = 0$ (rather than $\pi$) holds in the positive-definite neighborhood of the dominant $p=0$ Clock vacuum ($\det M_u \det M_d > 0$).
    \item \textbf{Non-Abelian Weyl--Heisenberg Completeness Preserves CKM CP Violation:} In an abelian $\Z_3$ flavor group (No-Go Theorem~\ref{thm:nogo2}), imposing $c_{-p,-q} = c_{p,q}^*$ forces $M_u$ and $M_d$ to be commuting circulants with $J_{\mathrm{CP}} = 0$. What our Weyl--Heisenberg reconstruction in $\R[\Delta(27)_f]$ proves is that the $9$ non-commuting operators $W_{p,q} = P^p Q^q \in \Delta(27)_f$ form a complete unitary operator basis ($\frac{1}{3}\Tr(W_{p,q}^\dagger W_{p',q'}) = \delta_{pp'}\delta_{qq'}$) spanning the \emph{entire} $9$-dimensional real space of $3 \times 3$ complex Hermitian matrices $H_3(\C)$, so a Hermitian $M_d = V_{\mathrm{CKM}} D_d V_{\mathrm{CKM}}^\dagger$ reproduces the full complex off-diagonal CKM phase ($\operatorname{Im}\Tr([M_u, M_d]^3) \neq 0$, $J_{\mathrm{CP}} = 2.93 \times 10^{-5}$) while keeping $\arg\det(M_u M_d) = 0$ at tree level.
    \item \textbf{Unprotected Radiative Renormalization (Open Problem):} Because $Q_L \sim \Sigma[18]$ and $q_R \sim \Sigma[10,11]$ have different gauge interactions and no left--right parity symmetry in $\Gstar$ exchanges them, the tree-level Hermiticity condition $c_{-p,-q} = c_{p,q}^*$ is not protected by a gauge symmetry under renormalization. If CP is violated spontaneously by flavon VEV phases, generic one-loop and two-loop diagrams involving the heavy 2HDM ($2.46\text{ TeV}$) and flavon ($5.17 \times 10^9\text{ GeV}$) scalars will generate non-Hermitian radiative corrections to $M_q$ and shift $\bar{\theta}_{\mathrm{QCD}}$ well above the pure-SM CKM-only estimate ($\Delta\bar{\theta}_{\mathrm{QCD}}^{\mathrm{SM}} \sim 10^{-16}$) unless suppressed by additional Nelson--Barr-type structure (Section~\ref{sec:open_problems}, item~3).
\end{enumerate}

\begin{thebibliography}{99}
\bibitem{Bombelli1987} L.~Bombelli, J.~Lee, D.~Meyer, and R.~D.~Sorkin, ``Space-time as a causal set,'' \emph{Phys.\ Rev.\ Lett.} \textbf{59}, 521 (1987).
\bibitem{ChamseddineConnes1997} A.~H.~Chamseddine and A.~Connes, ``The Spectral Action Principle,'' \emph{Commun.\ Math.\ Phys.} \textbf{186}, 731--750 (1997).
\bibitem{ShephardTodd1954} G.~C.~Shephard and J.~A.~Todd, ``Finite unitary reflection groups,'' \emph{Canad.\ J.\ Math.} \textbf{6}, 274--304 (1954).
\bibitem{Coxeter1974} H.~S.~M.~Coxeter, \emph{Regular Complex Polytopes} (Cambridge University Press, Cambridge, 1974).
\bibitem{MoodyPatera1993} R.~V.~Moody and J.~Patera, ``Quasicrystals and icosians,'' \emph{J.\ Phys.\ A: Math.\ Gen.} \textbf{26}, 2829--2853 (1993).
\bibitem{DistlerGaribaldi2010} J.~Distler and S.~Garibaldi, ``There is no `Theory of Everything' inside $E_8$,'' \emph{Commun.\ Math.\ Phys.} \textbf{298}, 419--436 (2010).
\bibitem{GeorgiJarlskog1979} H.~Georgi and C.~Jarlskog, ``A new lepton-quark mass relation in a unified theory,'' \emph{Phys.\ Lett.\ B} \textbf{86}, 297--300 (1979).
\bibitem{Koide1983} Y.~Koide, ``New view of quark and lepton mass hierarchy,'' \emph{Phys.\ Rev.\ D} \textbf{28}, 252 (1983).
\bibitem{IbanezRoss1991} L.~E.~Ib\'a\~nez and G.~G.~Ross, ``Discrete gauge symmetry anomalies,'' \emph{Phys.\ Lett.\ B} \textbf{260}, 291--295 (1991).
\bibitem{Cirelli2006} M.~Cirelli, N.~Fornengo, and A.~Strumia, ``Minimal dark matter,'' \emph{Nucl.\ Phys.\ B} \textbf{753}, 178--194 (2006).
\bibitem{Buttazzo2013} D.~Buttazzo \emph{et al.}, ``Investigating the near-criticality of the Higgs boson,'' \emph{JHEP} \textbf{12}, 089 (2013).
\bibitem{KleitmanRothschild1975} D.~J.~Kleitman and B.~L.~Rothschild, ``Asymptotic enumeration of partial orders on a finite set,'' \emph{Trans.\ Amer.\ Math.\ Soc.} \textbf{205}, 205--220 (1975).
\bibitem{BombelliHensonSorkin2009} L.~Bombelli, J.~Henson, and R.~D.~Sorkin, ``Discreteness without symmetry breaking: a theorem,'' \emph{Mod.\ Phys.\ Lett.\ A} \textbf{24}, 2579--2587 (2009).
\bibitem{GattoSartoriTonin1968} R.~Gatto, G.~Sartori, and M.~Tonin, ``Weak self-masses, Cabibbo angle, and broken $SU(2) \times SU(2)$,'' \emph{Phys.\ Lett.\ B} \textbf{28}, 128--130 (1968).
\bibitem{Georgi1979} H.~Georgi, ``Towards a grand unified theory of flavor,'' \emph{Nucl.\ Phys.\ B} \textbf{156}, 126--134 (1979).
\bibitem{Zeldovich1974} Ya.~B.~Zeldovich, I.~Yu.~Kobzarev, and L.~B.~Okun, ``Cosmological consequences of the spontaneous breakdown of discrete symmetry,'' \emph{Zh.\ Eksp.\ Teor.\ Fiz.} \textbf{67}, 3--11 (1974) [\emph{Sov.\ Phys.\ JETP} \textbf{40}, 1--5 (1975)].
\bibitem{PDG2024} S.~Navas \emph{et al.} (Particle Data Group), ``Review of Particle Physics,'' \emph{Phys.\ Rev.\ D} \textbf{110}, 030001 (2024).
\end{thebibliography}

\end{document}
